%=======================================================================
\chapter{Elastic Solids}
%=======================================================================

An elastic material is what is left of the general model of Chapter 8
when the constitutive functions are made blind to the velocity gradient.
Chapter 9 removed the reference configuration instead and arrived at
fluids; here the reference configuration stays and the rate of
deformation goes. The two deletions are independent, and the chapter
shows how much survives each one.

The list of response functions is that of Chapter 8,
\begin{equation}
   \mathcal{C}=\{\bT,\bq,\psi,\eta\},
   \tag{11.1}\label{eq:11.1}
\end{equation}
and the list of variables on which they are allowed to depend is
\eqr{8.4} with $\bL$ struck out,
\begin{equation}
   \mathcal{L}=\{\bx,\bv,\theta,\bF,\grad\theta\},
   \tag{11.2}\label{eq:11.2}
\end{equation}
both evaluated at the same material point $p$. For the Helmholtz free
energy of \eqr{6.186} this reads
\begin{equation}
   \psi=\hat\psi_{\kappa}(\bx,\bv,\theta,\bF_{\kappa},\grad\theta;p),
   \tag{11.3}\label{eq:11.3}
\end{equation}
the subscript $\kappa$ recording that the observer $O$ computes
$\bF_{\kappa}$ from a reference configuration of its own choosing.

\begin{remark}[what has been deleted, and what that costs]%
\label{rem:ch11delete}
Deleting $\bL$ is not a small change. In Chapter 9 the stress of a
viscous fluid was built out of $\bD$, and \eqr{9.38} showed that when
$\bD$ is switched off nothing remains but a pressure. Here $\bD$ is
switched off from the start, and yet Section~\ref{sec:ch11stress} will
produce a stress that depends on the deformation in full. The
difference is the reference configuration: a fluid has none to depend
on, so removing the rate leaves it with only $J$, while a solid
remembers $\kappa$ and can measure how far it has been carried from it.

The deletion also has a thermodynamical consequence that is not
obvious. For fluids, Chapter 9 had to \emph{assume} that the stress is
insensitive to $\grad\theta$. Here that insensitivity will be
\emph{proved}, at \eqr{11.16}$_{2}$, and the proof is the reason
\eqr{11.13} is worth writing out in full.
\end{remark}

Frame invariance, in the form \eqr{7.37} carried through the definition
\eqr{6.186} of $\psi$, requires that the constitutive function used by a
second observer $O^{+}$ deliver the same value:
\begin{equation}
   \hat\psi_{\kappa}(\bx,\bv,\theta,\bF_{\kappa},\grad\theta;p)
   =\hat\psi^{+}_{\kappa^{+}}\big(\bx^{+},\bv^{+},\theta^{+},
   \bF^{+}_{\kappa^{+}},\grad^{+}\theta^{+};p\big).
   \tag{11.4}\label{eq:11.4}
\end{equation}
Now run the argument of Sections~\ref{sec:ch9mfi} and
\ref{sec:ch10obs}: compare two relative motions of the observers, as at
\eqr{9.6}, chosen so that $\bQ_{2}(t_{0})=\bQ_{1}(t_{0})$ and
$\bc_{2}(t_{0})=\bc_{1}(t_{0})$, so that the two agree on the
configuration at $t_{0}$ and differ only in how they are moving. The
conclusion is
\begin{equation}
   \hat\psi_{\kappa}\big(\bx_{2},\bv_{2},\theta,\bF_{2},
   (\grad\theta)_{2};p\big)
   =\hat\psi_{\kappa}\big(\bx_{1},\bv_{1},\theta,\bF_{1},
   (\grad\theta)_{1};p\big),
   \tag{11.5}\label{eq:11.5}
\end{equation}
where $\bx_{2}$ and $\bv_{2}$ follow from $\bx_{1}$ and $\bv_{1}$ by
\eqr{9.12} with $\bQ=\bQ^{t}_{2}\bQ_{1}$ a rotation, and where
\begin{equation}
   \bF_{2}=\bQ\bF_{1}\qquad\text{and}\qquad
   (\grad\theta)_{2}=\bQ(\grad\theta)_{1}.
   \tag{11.6}\label{eq:11.6}
\end{equation}
Choosing $\bQ=\I$ leaves $\bF$ and $\grad\theta$ alone while moving
$\bx$ and $\bv$, so $\hat\psi_{\kappa}$ can depend on neither of the
latter two:
\begin{equation}
   \psi=\Psi(\theta,\bF,\grad\theta;p)
   \tag{11.7}\label{eq:11.7}
\end{equation}
for some function $\Psi$, the subscript $\kappa$ now suppressed.
Restoring general $\bQ$ in \eqr{11.5} then leaves
\begin{equation}
   \Psi(\theta,\bQ\bF,\bQ\grad\theta;p)=\Psi(\theta,\bF,\grad\theta;p)
   \qquad\text{for all rotations }\bQ .
   \tag{11.8}\label{eq:11.8}
\end{equation}
The identical reasoning applied to the stress and the heat flux removes
$\bx$ and $\bv$ from them too. For the heat flux,
\begin{equation}
   \bq=\hat\bq(\theta,\bF,\grad\theta;p),
   \tag{11.9}\label{eq:11.9}
\end{equation}
where, $\bq$ being a vector rather than a scalar, the transformation
carries a $\bQ$ on the left as well:
\begin{equation}
   \hat\bq(\theta,\bQ\bF,\bQ\grad\theta;p)
   =\bQ\,\hat\bq(\theta,\bF,\grad\theta;p)
   \qquad\text{for all rotations }\bQ .
   \tag{11.10}\label{eq:11.10}
\end{equation}

%=======================================================================
\section{Clausius--Duhem inequality}\label{sec:ch11cd}
%=======================================================================

Before exploiting \eqr{11.8} and \eqr{11.10}, which are statements about
symmetry, we spend the second law. It is \eqr{6.187},
\[
   -\rho\big(\dot\psi+\eta\dot\theta\big)+\bT\cdot\bL
   -\theta^{-1}\ip\bq{\grad\theta}\geq0 ,
\]
and the plan is to convert every term into a referential one and then
read off what the inequality forbids.

Two ingredients are needed. The first is the chain rule applied to
\eqr{11.7}, the free energy now being a function of three arguments:
\begin{equation}
   \dot\psi=\Psi_{\theta}\dot\theta+\Psi_{\bF}\cdot\dot\bF
   +\Psi_{\grad\theta}\cdot(\grad\theta)^{\textstyle\cdot} ,
   \tag{11.11}\label{eq:11.11}
\end{equation}
in which $\Psi_{\theta}$ is a scalar, $\Psi_{\bF}$ a two-point tensor
and $\Psi_{\grad\theta}$ a spatial vector, each paired with its argument
by the appropriate inner product. The second converts the stress power
density into one written on $\dot\bF$:
\begin{equation}
   \bT\cdot\bL=\bT\cdot\dot\bF\bF^{-1}
   =\tr\big(\bT\bF^{-t}\dot\bF^{t}\big)=\bT\bF^{-t}\cdot\dot\bF .
   \tag{11.12}\label{eq:11.12}
\end{equation}
The middle step is $\bA\cdot\bB=\tr(\bA\bB^{t})$ from
Definition~\ref{def:tip} with $\bB=\dot\bF\bF^{-1}$, whose transpose is
$\bF^{-t}\dot\bF^{t}$; the last step reads the same trace backwards with
the factors regrouped as $(\bT\bF^{-t})(\dot\bF)^{t}$.

Now multiply the inequality by $J$, which is positive and so does not
turn it around. Three conversions occur at once. By \eqr{2.59},
$J\rho=\rho_{\kappa}$. By \eqr{2.55} and \eqr{6.140},
$J\bT\bF^{-t}=\bT\bF^{*}=\bP$, so the stress power becomes
$\bP\cdot\dot\bF$, which is \eqr{6.156} again. The dissipation term
simply acquires the factor $J$. Substituting \eqr{11.11} and collecting
the coefficient of each rate,
\begin{equation}
\begin{aligned}
   -\rho_{\kappa}\big(\Psi_{\theta}+\eta\big)\dot\theta
   &+\big(\bP-\rho_{\kappa}\Psi_{\bF}\big)\cdot\dot\bF
   -\rho_{\kappa}\Psi_{\grad\theta}\cdot(\grad\theta)^{\textstyle\cdot}\\
   &-J\theta^{-1}\ip\bq{\grad\theta}\geq0 .
\end{aligned}
   \tag{11.13}\label{eq:11.13}
\end{equation}

This has the shape met at \eqr{9.64}. Write
\begin{equation}
\begin{aligned}
   \bz&=\big(\theta,\bF,\grad\theta\big),\qquad
   \bw=\big(\dot\theta,\dot\bF,(\grad\theta)^{\textstyle\cdot}\big),\\
   \ba(\bz)&=\big(-\rho_{\kappa}(\Psi_{\theta}+\eta),\
   \bP-\rho_{\kappa}\Psi_{\bF},\ -\rho_{\kappa}\Psi_{\grad\theta}\big)
\end{aligned}
   \tag{11.14}\label{eq:11.14}
\end{equation}
and
\begin{equation}
   b(\bz)=-J\theta^{-1}\ip\bq{\grad\theta},
   \tag{11.15}\label{eq:11.15}
\end{equation}
so that \eqr{11.13} reads $\ba(\bz)\cdot\bw+b(\bz)\geq0$. The
Coleman--Noll argument of \eqr{9.68} applies verbatim: at a fixed state
$\bz$ the rates $\bw$ may be assigned arbitrarily and independently, so
a non-zero $\ba(\bz)$ could be defeated by choosing $\bw=-\lambda\ba$
with $\lambda$ large. Hence $\ba(\bz)=\bze$ and $b(\bz)\geq0$. Written
out, and using $\Psi_{\grad\theta}=\bze$ to drop $\grad\theta$ from the
list,
\begin{equation}
   \eta=-\Psi_{\theta}(\theta,\bF;p),\qquad
   \bP=W_{\bF}(\theta,\bF;\bX),\qquad
   \ip\bq{\grad\theta}\leq0 ,
   \tag{11.16}\label{eq:11.16}
\end{equation}
where
\begin{equation}
   W(\theta,\bF;\bX)=\rho_{\kappa}(\bX)\,
   \Psi\big(\theta,\bF;\kappa^{-1}(\bX)\big)
   \tag{11.17}\label{eq:11.17}
\end{equation}
is the free energy per unit volume of $\kappa$. Passing from
$\rho_{\kappa}\Psi_{\bF}$ to $W_{\bF}$ is legitimate because
$\rho_{\kappa}$ depends on $\bX$ alone, by \eqr{6.18}, and so rides
through the differentiation with respect to $\bF$ as a constant. The
third of \eqr{11.16} follows from $b(\bz)\geq0$ after dividing by the
positive numbers $J$ and $\theta^{-1}$ and reversing the sign.

\begin{remark}[three results, and the one that is a surprise]%
\label{rem:ch11three}
The first of \eqr{11.16} is familiar: entropy is minus the temperature
derivative of the free energy, exactly as in equilibrium thermostatics.
The third is the statement that heat does not flow up a temperature
gradient, which is the local form of the second law one expects.

The second is the one worth pausing on, and it says two things at once.
That $\bP=W_{\bF}$ is the definition of hyperelasticity: the stress is
the gradient of a potential, so no closed cycle of deformation at fixed
temperature can extract work. Chapter 9 had no such statement, because
the viscous stress there is not a gradient and dissipates.

The other thing it says is hidden in the argument list. Nothing was
assumed about the dependence of $\bP$ on $\grad\theta$; what was assumed
is that $\Psi$ may depend on it, and the inequality then forced
$\Psi_{\grad\theta}=\bze$. The stress inherits that, and so the Piola
stress, and with it the Cauchy stress by \eqr{6.140}, cannot depend on
the temperature gradient at all. Remark~\ref{rem:ch11delete} promised
this. For viscous fluids the same insensitivity had to be put in by
hand, as Chapter 9 records.
\end{remark}

%=======================================================================
\section{Stress response}\label{sec:ch11stress}
%=======================================================================

With $\grad\theta$ gone from $\Psi$, the invariance requirement
\eqr{11.8} becomes a statement about $W$ alone. Multiplying \eqr{11.8}
by $\rho_{\kappa}$ and using \eqr{11.17},
\begin{equation}
   W(\theta,\bQ\bF;\bX)=W(\theta,\bF;\bX)
   \qquad\text{for all rotations }\bQ .
   \tag{11.18}\label{eq:11.18}
\end{equation}

\begin{proposition}[the content of \eqr{11.18}]\label{prop:ch11WC}
A function $W$ satisfies \eqr{11.18} if and only if it depends on $\bF$
only through $\bC=\bF^{t}\bF$, that is, if and only if
\begin{equation}
   W(\theta,\bF;\bX)=\bar W(\theta,\bC;\bX)
   \tag{11.19}\label{eq:11.19}
\end{equation}
for some function $\bar W$.
\end{proposition}

\begin{proof}
If \eqr{11.19} holds then \eqr{11.18} follows, because $\bQ\bF$ produces
the same $\bC$:
\[
   (\bQ\bF)^{t}(\bQ\bF)=\bF^{t}\bQ^{t}\bQ\bF=\bF^{t}\I\bF=\bC ,
\]
$\bQ$ being orthogonal.

Conversely, suppose \eqr{11.18}. Let $\bF=\bR_{F}\bU$ be the polar
decomposition \eqr{3.32}, and take
\[
   \bQ=\shift\bR^{t}_{F} .
\]
This is a rotation of $\Vt$ onto itself: $\bR^{t}_{F}$ carries $\Vt$ to
$\hVt$ and the shifter $\shift$ carries $\hVt$ back to $\Vt$, both
preserve inner products, and both have determinant $+1$. Substituting,
\[
   \bQ\bF=\shift\bR^{t}_{F}\bR_{F}\bU=\shift\hI\bU=\shift\bU ,
\]
so \eqr{11.18} gives $W(\theta,\bF;\bX)=W(\theta,\shift\bU;\bX)$. The
right-hand side involves $\bF$ only through $\bU$, and $\bU$ is
determined by $\bC$, being its unique positive-definite square root by
Theorem~\ref{thm:sqrt}. Defining $\bar W(\theta,\bC;\bX)$ to be
$W(\theta,\shift\bC^{1/2};\bX)$ gives \eqr{11.19}.
\end{proof}

\begin{remark}[why $\bC$ and not $\bU$]\label{rem:ch11whyC}
The proof produced $\bU$ and then traded it for $\bC$. Both are legal
arguments, and $\bC$ is chosen because it is a polynomial in $\bF$ while
$\bU$ is a square root: differentiating $\bar W$ with respect to $\bC$
is elementary, while differentiating with respect to $\bU$ drags in the
derivative of the square root, which is awkward: Theorem~\ref{thm:sqrt}
delivers $\bU$ but not a formula for it, and the closed form of
Theorem~\ref{thm:Uclosed} degenerates when two principal stretches
coincide. The same reason
made $\bC$ rather than $\bU$ the working strain measure at \eqr{2.97}.
\end{remark}

To convert \eqr{11.19} into a formula for the stress, differentiate it
along a motion. By the chain rule on both sides,
\begin{equation}
   W_{\bF}\cdot\dot\bF=\bar W_{\bC}\cdot\dot\bC
   =\bar W_{\bC}\cdot\big(\dot\bF^{t}\bF+\bF^{t}\dot\bF\big),
   \tag{11.20}\label{eq:11.20}
\end{equation}
where $\bar W_{\bC}=(\bar W_{\bC})^{t}$ is symmetric, being the gradient
of a scalar with respect to a symmetric tensor. The two terms are now
moved onto $\dot\bF$ using the rules
\begin{equation}
   \bA_{1}\cdot\bA_{2}\bA_{3}=\bA^{t}_{2}\bA_{1}\cdot\bA_{3}
   =\bA_{3}\bA^{t}_{1}\cdot\bA^{t}_{2},
   \tag{11.21}\label{eq:11.21}
\end{equation}
which are Problem 1 and follow from the cyclic property of the trace.
Applying the first of them to $\bar W_{\bC}\cdot\bF^{t}\dot\bF$, with
$\bA_{2}=\bF^{t}$, gives $\bF\bar W_{\bC}\cdot\dot\bF$. Applying the
second to $\bar W_{\bC}\cdot\dot\bF^{t}\bF$, with $\bA_{2}=\dot\bF^{t}$
and $\bA_{3}=\bF$, gives $\bF\bar W^{t}_{\bC}\cdot\dot\bF$, and the
symmetry of $\bar W_{\bC}$ makes this the same thing. The two terms are
therefore equal, and
\begin{equation}
   \big[W_{\bF}-2\bF\big(\bar W_{\bC}\big)\big]\cdot\dot\bF=0 .
   \tag{11.22}\label{eq:11.22}
\end{equation}
The bracket is a two-point tensor and $\dot\bF$ is an arbitrary
two-point tensor, the motion being at our disposal. Two-point tensors
form a linear space, so a fixed element of it whose inner product with
every element vanishes is itself zero. Hence
\begin{equation}
   W_{\bF}=2\bF\big(\bar W_{\bC}\big).
   \tag{11.23}\label{eq:11.23}
\end{equation}
Combining this with $\bP=W_{\bF}$ from \eqr{11.16}$_{2}$ and with
$\bP=\bF\bS$ from \eqr{6.144}, and cancelling the invertible $\bF$,
\begin{equation}
   \bS=2\bar W_{\bC} .
   \tag{11.24}\label{eq:11.24}
\end{equation}
This is automatically symmetric, and \eqr{6.145} then makes the Cauchy
stress symmetric as well.

\begin{remark}[rotational momentum balance comes for free]%
\label{rem:ch11free}
Chapter 6 obtained $\bT=\bT^{t}$ at \eqr{6.117} by localizing Euler's
second postulate, and Proposition~\ref{prop:ch6indep} showed that it is
genuinely independent information. Here it costs nothing: \eqr{11.24}
delivers a symmetric $\bS$ because $\bar W_{\bC}$ is the gradient of a
scalar with respect to a symmetric argument, and \eqr{6.145} transports
the symmetry to $\bT$.

Nothing has been proved twice. What happened is that \eqr{11.19}, which
came from frame invariance, already contains the symmetry: a constitutive
function of $\bC$ cannot produce an unsymmetric stress. Constitutive
theory and the balance laws have to be consistent, and this is the
consistency being checked. Had \eqr{11.24} come out unsymmetric, the
model would have been discarded.
\end{remark}

%-----------------------------------------------------------------------
\subsection{Material symmetry}\label{sec:ch11symm}
%-----------------------------------------------------------------------

Frame invariance restricts how a constitutive function responds to
rotating the \emph{observer}. Material symmetry restricts how it
responds to rotating the \emph{reference configuration}, and the two are
independent. In the form \eqr{8.16} and \eqr{8.17} it reads
\begin{equation}
   W(\theta,\bF;\bX)=W(\theta,\bF\bR;\bX)
   \qquad\text{for all }\bR\in\mathcal{G}_{\kappa}(p),
   \tag{11.25}\label{eq:11.25}
\end{equation}
$\mathcal{G}_{\kappa}(p)$ being the symmetry group at $p$ in the
configuration $\kappa$. Note where $\bR$ sits: on the right of $\bF$,
because it acts before the deformation, whereas the $\bQ$ of \eqr{11.18}
sits on the left and acts after it. Passing to $\bar W$ through
\eqr{11.19}, and using $(\bF\bR)^{t}(\bF\bR)=\bR^{t}\bC\bR$,
\begin{equation}
   \bar W(\theta,\bC;\bX)=\bar W(\theta,\bR^{t}\bC\bR;\bX).
   \tag{11.26}\label{eq:11.26}
\end{equation}

\begin{remark}[two rotations that must not be confused]%
\label{rem:ch11tworot}
The $\bR$ of \eqr{11.26} is a referential tensor, mapping $\hVt$ to
itself, and it is a member of the symmetry group. The $\bR_{F}$ of the
polar decomposition is a two-point tensor, mapping $\hVt$ to $\Vt$, and
it is determined by the deformation. They are different objects wearing
the same letter, and the book warns about it at this point. In these
notes the polar factor carries the subscript $F$ wherever both appear in
one argument, as in the proof of Proposition~\ref{prop:ch11WC} and in
\eqr{11.44} below.
\end{remark}

\subsection*{The unimodular case: fluids again}

Suppose $\mathcal{G}_{\kappa}(p)=U$, the unimodular group of
\eqr{8.18}. The argument of Chapter 9 then applies unchanged and forces
$W$ to depend on $\bF$ only through $J=\det\bF$, say
$W(\theta,\bF;\bX)=w(\theta,J;\bX)$. Hold $\theta$ fixed and
differentiate along a motion, using \eqr{11.16}$_{2}$ on the left and
Problem 2(b) of Chapter 4, $\partial J/\partial F_{iA}=F^{*}_{iA}$, on
the right:
\begin{equation}
   \bP\cdot\dot\bF=\dot W=\dot w=w_{J}\dot J
   =w_{J}\,\bF^{*}\cdot\dot\bF .
   \tag{11.27}\label{eq:11.27}
\end{equation}
As $\dot\bF$ is arbitrary this gives $\bP=w_{J}\bF^{*}$, and \eqr{6.140}
in the form $\bP=\bT\bF^{*}$ then yields, after cancelling the
invertible $\bF^{*}$,
\[
   \bT=w_{J}\,\I .
\]
Combined with conservation of mass, which ties $J$ to $\rho$ through
\eqr{2.59}, this is exactly the constitutive equation \eqr{9.38} of an
inviscid fluid, with the pressure $\tilde p=-w_{J}$. So the fluid is not
a separate theory: it is the elastic solid whose symmetry group is as
large as it can be.

\subsection*{Isotropy}

For a solid there is a local undistorted configuration in which the
symmetry group is contained in $\Orth^{+}$. If it is the whole of
$\Orth^{+}$ the material is called \emph{hemitropic} at the point: no
rotation applied before the deformation can be detected by any
experiment. Taking $\kappa(p)$ to be such a configuration,
\begin{equation}
   W(\theta,\bF;\bX)=W(\theta,\bF\bR;\bX)
   \qquad\text{for all }\bR\in\Orth^{+},
   \tag{11.28}\label{eq:11.28}
\end{equation}
and therefore, by \eqr{11.26},
\begin{equation}
   \bar W(\theta,\bC;\bX)=\bar W(\theta,\bR^{t}\bC\bR;\bX)
   \qquad\text{for all }\bR\in\Orth ,
   \tag{11.29}\label{eq:11.29}
\end{equation}
now over the \emph{full} orthogonal group. The upgrade is free because
$\Orth=\Orth^{+}\cup\{-\hI\}\Orth^{+}$: every improper orthogonal tensor
is $-\bR$ for some rotation $\bR$, and
\[
   (-\bR)^{t}\bC(-\bR)=(-1)^{2}\bR^{t}\bC\bR=\bR^{t}\bC\bR ,
\]
so the two signs give the same restriction. A function obeying
\eqr{11.29} is called \emph{isotropic}.

The representation theorem proved in Part 1 of
Section~\ref{sec:ch9isotropic} now applies, a scalar isotropic function
of a symmetric tensor being a function of its principal invariants:
\begin{equation}
   \bar W(\theta,\bC;\bX)
   =\hat W\big(\theta,I_{1}(\bC),I_{2}(\bC),I_{3}(\bC);\bX\big),
   \tag{11.30}\label{eq:11.30}
\end{equation}
where
\begin{equation}
   I_{1}(\bC)=\tr\bC,\qquad
   I_{2}(\bC)=\tfrac12\big[I^{2}_{1}-\tr(\bC^{2})\big],\qquad
   I_{3}(\bC)=\det\bC
   \tag{11.31}\label{eq:11.31}
\end{equation}
are the principal invariants of \eqr{3.2}. Substituting into
\eqr{11.24} and writing $\hat W_{k}=\partial\hat W/\partial I_{k}$,
\begin{equation}
   \Big[\tfrac12\bS-\sum^{3}_{k=1}\hat W_{k}\,(I_{k})_{\bC}\Big]
   \cdot\dot\bC=0\qquad\text{for all symmetric }\dot\bC ,
   \tag{11.32}\label{eq:11.32}
\end{equation}
and the gradients of the invariants, computed in Section B.3 of
Appendix B and reproduced here for use,
\[
   (I_{1})_{\bC}=\hI,\qquad
   (I_{2})_{\bC}=I_{1}\hI-\bC,\qquad
   (I_{3})_{\bC}=I_{3}\bC^{-1},
\]
deliver the Piola--Kirchhoff stress
\begin{equation}
   \bS=2\big[\big(\hat W_{1}+I_{1}\hat W_{2}\big)\hI
   -\hat W_{2}\bC+I_{3}\hat W_{3}\bC^{-1}\big].
   \tag{11.33}\label{eq:11.33}
\end{equation}
Pushing this forward with \eqr{6.145}, $J\bT=\bF\bS\bF^{t}$, and using
\[
\begin{aligned}
   \bF\hI\bF^{t}&=\bB,\qquad
   \bF\bC\bF^{t}=\bF\bF^{t}\bF\bF^{t}=\bB^{2},\\
   \bF\bC^{-1}\bF^{t}&=\bF\bF^{-1}\bF^{-t}\bF^{t}=\I ,
\end{aligned}
\]
the middle one by $\bC=\bF^{t}\bF$ and the last by
$\bC^{-1}=\bF^{-1}\bF^{-t}$, gives the Cauchy stress
\begin{equation}
   J\bT=2\big[\big(\hat W_{1}+I_{1}\hat W_{2}\big)\bB
   -\hat W_{2}\bB^{2}+I_{3}\hat W_{3}\I\big],
   \tag{11.34}\label{eq:11.34}
\end{equation}
$\bB$ being the left Cauchy--Green tensor.

The invariants appearing here may be read off either tensor, because
\begin{equation}
\begin{aligned}
   I_{1}(\bC)&=\lambda^{2}_{1}+\lambda^{2}_{2}+\lambda^{2}_{3}
   =I_{1}(\bB),\\
   I_{2}(\bC)&=\lambda^{2}_{1}\lambda^{2}_{2}
   +\lambda^{2}_{2}\lambda^{2}_{3}+\lambda^{2}_{1}\lambda^{2}_{3}
   =I_{2}(\bB),\\
   I_{3}(\bC)&=\lambda^{2}_{1}\lambda^{2}_{2}\lambda^{2}_{3}=I_{3}(\bB),
\end{aligned}
   \tag{11.35}\label{eq:11.35}
\end{equation}
the $\lambda_{k}$ being the principal stretches. This is Problem 3, and
it follows from the spectral representations \eqr{3.38}: $\bC$ and $\bB$
carry the same eigenvalues $\lambda^{2}_{k}$ on different eigenvectors,
and the invariants see only the eigenvalues.

If the material is incompressible then $J=1$, so $I_{3}=1$ throughout
the motion and $\hat W_{3}$ is never sampled. The term it multiplied is
replaced by the reaction of Chapter 10:
\begin{equation}
\begin{aligned}
   \bT&=-p\I+2\big[\big(W^{*}_{1}+I_{1}W^{*}_{2}\big)\bB
   -W^{*}_{2}\bB^{2}\big],\\
   &\text{where}\quad W^{*}(I_{1},I_{2};\bX)=\hat W(I_{1},I_{2},1;\bX),
\end{aligned}
   \tag{11.36}\label{eq:11.36}
\end{equation}
with $W^{*}_{k}=\partial W^{*}/\partial I_{k}$ and $p$ a constitutively
indeterminate pressure, the multiplier of
Remark~\ref{rem:ch10lagrange}.

\begin{remark}[which symmetries a constraint tolerates]%
\label{rem:ch11constraint}
The constraint \eqr{10.11} of incompressibility, $\det\bC=1$, is
unchanged when $\bC$ is replaced by $\bR^{t}\bC\bR$ for any
$\bR\in\Orth$, since $\det(\bR^{t}\bC\bR)=\det\bC$. So isotropy and
incompressibility can hold together. The same computation works for any
$\bR\in U$, the determinant being all a unimodular map preserves, so
incompressibility is compatible with every symmetry group, fluidity
included.

Inextensibility is different. The constraint \eqr{10.12} singles out a
direction $\bM$, and there are infinitely many $\bR\in\Orth$ that move
it. A material reinforced by a family of inextensible fibres therefore
cannot be isotropic, which is Problem 5 and which
Example~\ref{ex:ch10fibre} already suggested: the fibre is visible, so
some rotations of the reference configuration are detectable.
\end{remark}

\subsection*{Example: simple shear}

Take the simple shear of Section~\ref{sec:shear} and let it be
suffered by a uniform isotropic material, so that $\hat W$, or $W^{*}$
in the incompressible case, carries no explicit $\bX$.

The deformation is homogeneous, so by \eqr{3.54} the left Cauchy--Green
tensor is a constant field. Hence the stress delivered by \eqr{11.34}
with $J=1$, or by \eqr{11.36}, is also constant, its divergence
vanishes, and \eqr{6.109} is satisfied with zero body force: simple
shear is an equilibrium deformation for every uniform isotropic
material. For the incompressible case one condition is added, that
$\dvg(p\I)=\grad p$ vanish, so the reactive pressure must be uniform.

From \eqr{3.54}, in the basis $\{\be_{i}\}$,
\[
   (B_{ij})=\begin{pmatrix}1+\gamma^{2}&\gamma&0\\
   \gamma&1&0\\ 0&0&1\end{pmatrix},
\]
and squaring the matrix, entry by entry,
\begin{equation}
\begin{aligned}
   \bB^{2}&=\big[(1+\gamma^{2})^{2}+\gamma^{2}\big]\,
     \be_{1}\otimes\be_{1}
    +(1+\gamma^{2})\,\be_{2}\otimes\be_{2}+\be_{3}\otimes\be_{3}\\
   &\quad+\gamma(2+\gamma^{2})\,
     \big(\be_{1}\otimes\be_{2}+\be_{2}\otimes\be_{1}\big),
\end{aligned}
   \tag{11.37}\label{eq:11.37}
\end{equation}
the $11$ entry being $(1+\gamma^{2})^{2}+\gamma\cdot\gamma$, the $12$
entry $(1+\gamma^{2})\gamma+\gamma\cdot1=\gamma(2+\gamma^{2})$, and the
$22$ entry $\gamma\cdot\gamma+1\cdot1=1+\gamma^{2}$.

\begin{remark}[a misprint in the printed \eqr{11.37} and
\eqr{11.39}]\label{rem:ch11mis37}
The text prints the $11$ component of $\bB^{2}$ as
$(1+\gamma^{2})+\gamma^{2}$, both in \eqr{11.37} and in the third matrix
of \eqr{11.39}. It should be $(1+\gamma^{2})^{2}+\gamma^{2}$. Squaring
the matrix above, the $11$ entry is
\[
   B_{11}B_{11}+B_{12}B_{21}+B_{13}B_{31}
   =(1+\gamma^{2})^{2}+\gamma^{2} ,
\]
and at $\gamma=1$ this is $5$, whereas the printed expression gives $3$;
direct squaring of $\left(\begin{smallmatrix}2&1&0\\1&1&0\\
0&0&1\end{smallmatrix}\right)$ confirms $5$.

The error is not confined to that entry. Section~\ref{sec:ch11normal}
below derives the universal relation \eqr{11.42} from these components,
and with the printed value it comes out as
$T_{11}-T_{22}=\gamma^{2}(f_{1}+f_{2})$, which is not $\gamma T_{12}$
unless $f_{2}=0$. With the corrected value it comes out exactly. So
\eqr{11.42}, which the text calls remarkable and which is the
experimentally interesting statement of the section, is inconsistent
with \eqr{11.37} as printed and consistent with it as corrected.
\end{remark}

The invariants follow. With $\gamma$ the uniform shear,
\begin{equation}
   I_{1}(\bB)=I_{2}(\bB)=3+\gamma^{2},
   \tag{11.38}\label{eq:11.38}
\end{equation}
and $I_{3}(\bB)=1$. The first is a trace, $(1+\gamma^{2})+1+1$. The
second uses \eqr{11.31}$_{2}$ with the trace of \eqr{11.37},
\[
   \tr\bB^{2}=(1+\gamma^{2})^{2}+\gamma^{2}+(1+\gamma^{2})+1
   =3+4\gamma^{2}+\gamma^{4},
\]
so that
\[
   I_{2}=\tfrac12\big[(3+\gamma^{2})^{2}-(3+4\gamma^{2}+\gamma^{4})\big]
   =\tfrac12\big[9+6\gamma^{2}+\gamma^{4}-3-4\gamma^{2}-\gamma^{4}\big]
   =3+\gamma^{2},
\]
and $I_{3}=\det\bB=(\det\bF)^{2}=1$. This is Problem 6. Every
coefficient in \eqr{11.34} and \eqr{11.36} is therefore a function of
$\gamma^{2}$ alone.

Resolving those equations in the basis $\{\be_{i}\otimes\be_{j}\}$ and
collecting components,
\begin{equation}
\begin{aligned}
   (T_{ij})&=f_{0}\begin{pmatrix}1&0&0\\0&1&0\\0&0&1\end{pmatrix}
   +f_{1}(\gamma^{2})\begin{pmatrix}1+\gamma^{2}&\gamma&0\\
   \gamma&1&0\\0&0&1\end{pmatrix}\\[2pt]
   &\quad+f_{2}(\gamma^{2})\begin{pmatrix}
   (1+\gamma^{2})^{2}+\gamma^{2}&\gamma(2+\gamma^{2})&0\\
   \gamma(2+\gamma^{2})&1+\gamma^{2}&0\\ 0&0&1\end{pmatrix},
\end{aligned}
   \tag{11.39}\label{eq:11.39}
\end{equation}
where $f_{1}$ and $f_{2}$ are constitutive functions and $f_{0}$ is
another one in the compressible case, or minus the uniform constraint
pressure in the incompressible case. Comparing with \eqr{11.34}, the
three matrices are those of $\I$, $\bB$ and $\bB^{2}$.

%-----------------------------------------------------------------------
\subsection*{The normal stress effect}\label{sec:ch11normal}
%-----------------------------------------------------------------------

Reading \eqr{11.39} off the third row and column, $T_{13}=T_{23}=0$:
the shear does not drag the $\be_{3}$ direction. The shear component is
\begin{equation}
   T_{12}=\gamma\mu(\gamma^{2}),\qquad\text{where}\quad
   \mu(\gamma^{2})=f_{1}(\gamma^{2})+(2+\gamma^{2})f_{2}(\gamma^{2}),
   \tag{11.40}\label{eq:11.40}
\end{equation}
obtained by adding the $12$ entries, $f_{1}\gamma+f_{2}\gamma(2+
\gamma^{2})$, and taking out the common factor $\gamma$.

\begin{remark}[a second misprint, in \eqr{11.40}]\label{rem:ch11mis40}
The text prints $\mu(\gamma^{2})=f_{1}(\gamma^{2})+\gamma(2+\gamma^{2})
f_{2}(\gamma^{2})$, carrying a factor $\gamma$ that does not belong.
Two independent checks reject it. First, the $12$ entry of \eqr{11.39}
is $f_{1}\gamma+f_{2}\gamma(2+\gamma^{2})$, and dividing by $\gamma$
leaves $f_{1}+(2+\gamma^{2})f_{2}$. Second, the notation
$\mu(\gamma^{2})$ asserts that $\mu$ depends on $\gamma$ only through
$\gamma^{2}$, and the printed expression contains a term odd in
$\gamma$, so it contradicts its own argument list. The same factor also
breaks \eqr{11.42}.
\end{remark}

If $\mu$ is differentiable then, since it is a function of $\gamma^{2}$,
its expansion about $\gamma=0$ has no linear term:
\begin{equation}
   T_{12}=G\gamma+O(\gamma^{3}),
   \tag{11.41}\label{eq:11.41}
\end{equation}
where $G=\mu(0)$ is the classical shear modulus of small deformations.
The correction is cubic, not quadratic, and that is a consequence of
isotropy: reversing the shear must reverse the shear stress.

The diagonal components $T_{11}$, $T_{22}$ and $T_{33}$ are in general
non-zero, and they satisfy
\begin{equation}
   T_{11}-T_{22}=\gamma\,T_{12} .
   \tag{11.42}\label{eq:11.42}
\end{equation}
To see it, subtract the $22$ entries of \eqr{11.39} from the $11$
entries. The $f_{0}$ terms cancel, being equal, and
\[
\begin{aligned}
   T_{11}-T_{22}&=f_{1}\big[(1+\gamma^{2})-1\big]
   +f_{2}\big[(1+\gamma^{2})^{2}+\gamma^{2}-(1+\gamma^{2})\big]\\
   &=f_{1}\gamma^{2}
   +f_{2}\big[(1+\gamma^{2})\gamma^{2}+\gamma^{2}\big]
   =\gamma^{2}\big[f_{1}+(2+\gamma^{2})f_{2}\big]
   =\gamma^{2}\mu ,
\end{aligned}
\]
where the middle step factored $(1+\gamma^{2})^{2}-(1+\gamma^{2})
=(1+\gamma^{2})\gamma^{2}$. Comparing with \eqr{11.40},
$\gamma T_{12}=\gamma^{2}\mu$ as well, which is \eqr{11.42}. This is
Problem 7.

\begin{remark}[what makes \eqr{11.42} worth having]%
\label{rem:ch11universal}
No constitutive function survives in \eqr{11.42}. Whatever $\hat W$ or
$W^{*}$ may be, and whether the material is compressible or not, a
uniform isotropic body in simple shear satisfies it. A relation of this
kind is called \emph{universal}, and it is the rarest and most valuable
thing a nonlinear theory can produce, because it can be checked in the
laboratory without first identifying the material.

What it predicts is that shearing a block requires normal tractions as
well as shear tractions, and that the difference between two of them is
fixed by the shear stress and the shear. Hold the block between two
plates and slide one of them: the block tries to change its thickness,
and $T_{11}-T_{22}$ is what must be applied to stop it. The right-hand
side is $\gamma T_{12}=\gamma^{2}\mu$, of order $\gamma^{2}$, so the
linear theory of Chapter 12, which discards everything beyond first
order in the displacement gradient, predicts none of it. The normal
stress effect is inherently nonlinear.

Pure simple shear is hard to realise in a laboratory. Torsion is not,
and Section~\ref{sec:torsion} showed that torsion is locally a simple
shear, which is how the effect has been confirmed: a twisted bar tries
to lengthen. Problem 12 works that case out.
\end{remark}

\begin{figure}[htbp]\centering
\renewcommand{\thefigure}{A}
\begin{tikzpicture}[scale=1.05,
   ar/.style={-{Stealth[length=2.3mm]},thick},
   thn/.style={-{Stealth[length=1.9mm]},semithick,gray!80}]
% ---------------- left: what the linear theory expects -------------
\begin{scope}[xshift=0cm]
  \draw[guide] (0,0) rectangle (2.0,2.0);
  \draw[thick] (0,0)--(2.0,0)--(2.7,2.0)--(0.7,2.0)--cycle;
  \foreach \x in {0.9,1.4,1.9,2.4}{\draw[thn] (\x,2.16)--++(0.42,0);}
  \foreach \x in {0.15,0.65,1.15,1.65}{\draw[thn] (\x+0.42,-0.16)--++(-0.42,0);}
  \node[font=\small,anchor=south] at (1.7,2.30) {$T_{12}$};
  \node[font=\footnotesize,text=gray!80,anchor=north,align=center]
     at (1.35,-0.46) {shear tractions only};
  \node[font=\small,anchor=north] at (1.35,-1.02) {$T_{11}=T_{22}$};
\end{scope}
% ---------------- right: what the nonlinear theory requires ---------
\begin{scope}[xshift=5.4cm]
  \draw[guide] (0,0) rectangle (2.0,2.0);
  \draw[thick] (0,0)--(2.0,0)--(2.7,2.0)--(0.7,2.0)--cycle;
  \foreach \x in {0.9,1.4,1.9,2.4}{\draw[thn] (\x,2.16)--++(0.42,0);}
  \foreach \x in {0.15,0.65,1.15,1.65}{\draw[thn] (\x+0.42,-0.16)--++(-0.42,0);}
  \node[font=\small,anchor=south] at (1.7,2.30) {$T_{12}$};
  % the normal tractions that must be supplied as well
  \draw[ar] (1.35,2.62)--(1.35,2.16);
  \draw[ar] (1.35,-0.62)--(1.35,-0.16);
  \node[font=\small,anchor=west] at (1.46,2.62) {$T_{22}$};
  \draw[ar] (-0.55,1.0)--(-0.09,1.0);
  \draw[ar] (3.25,1.0)--(2.79,1.0);
  \node[font=\small,anchor=east] at (-0.62,1.0) {$T_{11}$};
  \node[font=\footnotesize,text=gray!80,anchor=north,align=center]
     at (1.35,-1.02) {normal tractions too};
  \node[font=\small,anchor=north] at (1.35,-1.52)
     {$T_{11}-T_{22}=\gamma T_{12}$};
\end{scope}
\end{tikzpicture}
\caption{The normal stress effect of \eqr{11.42}, with the undeformed
block dashed. Sliding the top face of a block over the bottom one is
not achieved by shear tractions alone. A uniform isotropic solid held
in simple shear also pushes on the faces normal to $\be_{1}$ and
$\be_{2}$, and the difference of those two normal stresses is
$\gamma T_{12}$, whatever the material. Left is what the linear theory
of Chapter 12 predicts, right what \eqr{11.34} and \eqr{11.36} require;
the difference is of order $\gamma^{2}$ and vanishes with the
approximation.}
\label{fig:ch11normal}
\end{figure}

%=======================================================================
\section{Heat flux}\label{sec:ch11heat}
%=======================================================================

The stress has been reduced by \eqr{11.19} and then by isotropy. The
same two steps are now applied to the heat flux, and the first of them
has a different flavour, because $\hat\bq$ carries a free vector index
and cannot be handled by the polar trick alone.

Let the material be hemitropic at $p$, so that the symmetry condition
on the constitutive function for the heat flux reads
\begin{equation}
   \hat\bq(\theta,\bF,\grad\theta;p)
   =\hat\bq(\theta,\bF\bR,\grad\theta;p)
   \qquad\text{for all }\bR\in\Orth^{+}.
   \tag{11.43}\label{eq:11.43}
\end{equation}
For a necessary condition, choose the particular member
$\bR=\bR^{t}_{F}\shift$, where $\shift$ is the shifter and $\bR_{F}$ the
two-point rotation of the polar factorization of $\bF$. It is a
referential rotation, being a composition of two orientation-preserving
isometries $\hVt\to\Vt\to\hVt$. Using $\bF=\bV\bR_{F}$ from \eqr{3.32},
\begin{equation}
   \hat\bq(\theta,\bF,\grad\theta;p)
   =\hat\bq(\theta,\bV\bR_{F}\bR^{t}_{F}\shift,\grad\theta;p)
   =\hat\bq(\theta,\bV\shift,\grad\theta;p),
   \tag{11.44}\label{eq:11.44}
\end{equation}
since $\bR_{F}\bR^{t}_{F}=\I$. The shifter is fixed in the frame of $O$
and carries no information about the deformation, and $\bV$ is
determined by the left Cauchy--Green tensor $\bB=\bF\bF^{t}=\bV^{2}$,
uniquely by Theorem~\ref{thm:sqrt}. So the right-hand side is a function
of $\bB$:
\begin{equation}
   \hat\bq(\theta,\bF,\grad\theta;p)=\bq^{*}(\theta,\bB,\grad\theta;\bX)
   \tag{11.45}\label{eq:11.45}
\end{equation}
for some function $\bq^{*}$. The condition is also sufficient.
Substituting \eqr{11.45} and using $\bR\bR^{t}=\hI$,
\begin{equation}
\begin{aligned}
   \hat\bq(\theta,\bF\bR,\grad\theta;p)
   &=\bq^{*}(\theta,\bF\bR\bR^{t}\bF^{t},\grad\theta;\bX)\\
   &=\bq^{*}(\theta,\bF\hI\bF^{t},\grad\theta;\bX)\\
   &=\bq^{*}(\theta,\bB,\grad\theta;\bX)
    =\hat\bq(\theta,\bF,\grad\theta;p)
\end{aligned}
   \tag{11.46}\label{eq:11.46}
\end{equation}
for every $\bR\in\Orth^{+}$, which is \eqr{11.43}. So \eqr{11.45} is
necessary and sufficient for hemitropy.

Frame invariance \eqr{11.10} may now be rewritten on $\bq^{*}$. Under
$\bF\mapsto\bQ\bF$ the left Cauchy--Green tensor becomes
$\bQ\bF\bF^{t}\bQ^{t}=\bQ\bB\bQ^{t}$, so
\begin{equation}
\begin{aligned}
   \bq^{*}(\theta,\bQ\bB\bQ^{t},\bQ\grad\theta;\bX)
   &=\bQ\,\bq^{*}(\theta,\bB,\grad\theta;\bX)\\
   &\text{for all rotations }\bQ\colon\Vt\to\Vt .
\end{aligned}
   \tag{11.47}\label{eq:11.47}
\end{equation}

\begin{remark}[why $\bQ$ is a rotation here and orthogonal in
Chapter 9]\label{rem:ch11whyrot}
The restriction \eqr{9.40} on the heat flux in a fluid looked the same
but ranged over all orthogonal $\bQ$. The difference is the reference
configuration. For a fluid the symmetry group is unimodular, so
Chapter 9 was able to absorb $\kappa$ entirely and the argument lost
track of orientation. For a solid $\kappa$ survives explicitly, and it
was chosen occupiable, so only orientation-preserving maps are
admissible; the discussion leading to \eqr{10.8} made the same point.
\end{remark}

The restriction to rotations can be lifted at the cost of one further
physical assumption. Suppose the heat flux is odd in the temperature
gradient, as it is in the Fourier law \eqr{9.61} and as it is in every
material for which measurements exist:
\begin{equation}
   \bq^{*}(\theta,\bB,-\grad\theta;\bX)
   =-\bq^{*}(\theta,\bB,\grad\theta;\bX).
   \tag{11.48}\label{eq:11.48}
\end{equation}
Then \eqr{11.47} holds for every orthogonal $\bQ$. For if $\bQ$ is
improper, write $\bQ=-\bQ'$ with $\bQ'=-\bQ$ a rotation, and note
$\bQ\bB\bQ^{t}=\bQ'\bB\bQ'^{t}$, the two signs cancelling, while
$\bQ\grad\theta=-\bQ'\grad\theta$. Hence
\[
\begin{aligned}
   \bq^{*}(\theta,\bQ\bB\bQ^{t},\bQ\grad\theta;\bX)
   &=\bq^{*}(\theta,\bQ'\bB\bQ'^{t},-\bQ'\grad\theta;\bX)\\
   &=-\bq^{*}(\theta,\bQ'\bB\bQ'^{t},\bQ'\grad\theta;\bX)\\
   &=-\bQ'\bq^{*}(\theta,\bB,\grad\theta;\bX)\\
   &=\bQ\,\bq^{*}(\theta,\bB,\grad\theta;\bX),
\end{aligned}
\]
the second equality by \eqr{11.48} and the third by \eqr{11.47} for the
rotation $\bQ'$.

With $\bQ$ now ranging over the full orthogonal group, the
representation \eqr{9.41} with \eqr{9.42} applies and gives
\begin{equation}
   \bq=\big(\phi_{0}\I+\phi_{1}\bB+\phi_{2}\bB^{2}\big)\grad\theta ,
   \tag{11.49}\label{eq:11.49}
\end{equation}
where $\phi_{0},\phi_{1},\phi_{2}$ are functions of $\theta$, of the
invariants $I_{1}(\bB),I_{2}(\bB),I_{3}(\bB)$, and of
\begin{equation}
\begin{aligned}
   I_{4}(\bB,\grad\theta)&=\ip{\grad\theta}{\bB\grad\theta},\qquad
   I_{5}(\bB,\grad\theta)=\ip{\grad\theta}{\bB^{2}\grad\theta},\\
   I_{6}(\grad\theta)&=|\grad\theta| ,
\end{aligned}
   \tag{11.50}\label{eq:11.50}
\end{equation}
and, for a non-uniform material, of $\bX$ as well.

%-----------------------------------------------------------------------
\subsection{General form of the heat flux}\label{sec:ch11heatgen}
%-----------------------------------------------------------------------

Equation \eqr{11.49} is isotropic, and it exhibits the heat flux as a
linear map applied to $\grad\theta$. That linearity is not a
consequence of isotropy. It holds for every symmetry group, and the
reason is the dissipation inequality \eqr{11.16}$_{3}$ alone.

\begin{remark}[a change of letter]\label{rem:ch11gclash}
The book calls the temperature-gradient variable $\bp$ in what follows.
That letter already carries two meanings in this chapter, the material
point $p$ of \eqr{11.3} and the reaction pressure $p$ of \eqr{11.36},
and the derivative $\nabla\bh(\bp)$ would then be a gradient with
respect to something that also indexes the body. These notes write
$\bg$ for the variable and keep $\bq$ for the flux.
\end{remark}

The claim to be established is
\begin{equation}
   \bq=\bK(\theta,\bF,\grad\theta;\bX)\,\grad\theta
   \tag{11.51}\label{eq:11.51}
\end{equation}
for a tensor-valued function $\bK$, a result of [22] holding for every
symmetry group.

Fix $\theta$, $\bF$ and $\bX$, and hence $p$, and define
\begin{equation}
   \bh(\bg)=\hat\bq(\theta,\bF,\bg;p),
   \tag{11.52}\label{eq:11.52}
\end{equation}
a vector-valued function of the single vector variable $\bg$. The
inequality \eqr{11.16}$_{3}$ says
\begin{equation}
   \ip\bg{\bh(\bg)}\leq0\qquad\text{for all }\bg .
   \tag{11.53}\label{eq:11.53}
\end{equation}

\begin{proposition}[no gradient, no flux]\label{prop:ch11noflux}
If $\bh$ is differentiable then $\bh(\bze)=\bze$; that is,
$\hat\bq(\theta,\bF,\bze;p)=\bze$.
\end{proposition}

\begin{proof}
The scalar function $\bg\mapsto\ip\bg{\bh(\bg)}$ is bounded above by
zero, by \eqr{11.53}, and it attains the value zero at $\bg=\bze$. So it
is maximized there, and being differentiable it is stationary there. Its
differential at $\bg$, in the direction $\dl\bg$, is
\begin{equation}
   \dl\ip\bg{\bh(\bg)}
   =\ip{\dl\bg}{\bh(\bg)}+\ip\bg{[\nabla\bh(\bg)]\dl\bg} ,
   \tag{11.54}\label{eq:11.54}
\end{equation}
$\nabla\bh$ being the tensor-valued gradient of $\bh$. Evaluating at
$\bg=\bze$ kills the second term, whatever $\nabla\bh(\bze)$ may be, and
stationarity leaves
\[
   \ip{\dl\bg}{\bh(\bze)}=0\qquad\text{for every }\dl\bg ,
\]
whence $\bh(\bze)=\bze$ by Lemma~\ref{lem:cvec}.
\end{proof}

\begin{equation}
   \hat\bq(\theta,\bF,\bze;p)=\bze .
   \tag{11.55}\label{eq:11.55}
\end{equation}

So there is no heat flux in the absence of a temperature gradient. The
statement is worth isolating, because it is often taken as a modelling
assumption and it is in fact a theorem of the second law.

Now fix $\bg$ and travel out to it along a straight line. Put
\begin{equation}
   \by(x)=\bh(x\bg),\qquad x\in[0,1] ,
   \tag{11.56}\label{eq:11.56}
\end{equation}
a curve in $\Vt$ joining $\by(0)=\bh(\bze)=\bze$ to $\by(1)=\bh(\bg)$.
By the chain rule, the argument $x\bg$ having derivative $\bg$ with
respect to $x$,
\begin{equation}
   \by'(x)=\big[\nabla\bh(x\bg)\big]\bg .
   \tag{11.57}\label{eq:11.57}
\end{equation}
Integrating from $0$ to $1$ and using $\by(0)=\bze$ from
Proposition~\ref{prop:ch11noflux},
\begin{equation}
   \bh(\bg)=\int^{1}_{0}\by'(x)\,\dl x=\big[\bK(\bg)\big]\bg ,
   \qquad\text{where}\quad
   \bK(\bg)=\int^{1}_{0}\nabla\bh(x\bg)\,\dl x ,
   \tag{11.58}\label{eq:11.58}
\end{equation}
the constant vector $\bg$ coming outside the integral because the
integration is over $x$. Reinstating the passive variables gives
\eqr{11.51}.

\begin{remark}[what \eqr{11.51} is and is not]\label{rem:ch11Kread}
It is not a linear constitutive law. The tensor $\bK$ depends on
$\grad\theta$, so \eqr{11.51} is a rewriting rather than a restriction:
it says only that the flux vanishes with the gradient and that the
dependence can always be factored through a tensor. Fourier's law is
the special case in which $\bK$ does not depend on $\grad\theta$, and
\eqr{11.49} is the special case in which $\bK$ is the isotropic
combination in parentheses.

Nor is $\bK$ unique or symmetric. The construction integrates
$\nabla\bh$ along one particular path, the straight line, and a
different path would give a different $\bK$ with the same product
$\bK\grad\theta$. In \eqr{11.49} the $\bK$ that appears is symmetric,
being a polynomial in the symmetric $\bB$; in general it is not.

That last point costs nothing, because the dissipation inequality only
ever sees the symmetric part. Writing \eqr{11.53} with \eqr{11.51},
\begin{equation}
   \bK\cdot\bg\otimes\bg\leq0 ,
   \tag{11.59}\label{eq:11.59}
\end{equation}
since $\bK\cdot(\bg\otimes\bg)=K_{ij}g_{i}g_{j}=\ip\bg{\bK\bg}$, and a
quadratic form is blind to the skew part of its matrix: $\ip\bg{\bK\bg}
=\ip\bg{(\Sym\bK)\bg}$ by Lemma~\ref{lem:symskw}, $\bg\otimes\bg$ being
symmetric.
\end{remark}

%=======================================================================
\setcounter{problemenv}{0}
\section{Problems}\label{sec:problems:c11}
%=======================================================================

\begin{problem}[The three rearrangement rules]
Prove \eqr{11.21}.
\end{problem}

\begin{solution}
All three expressions are the same trace, read with the factors grouped
differently. Start from Definition~\ref{def:tip},
$\bA\cdot\bB=\tr(\bA\bB^{t})$, and expand the left member:
\[
   \bA_{1}\cdot\bA_{2}\bA_{3}=\tr\big(\bA_{1}(\bA_{2}\bA_{3})^{t}\big)
   =\tr\big(\bA_{1}\bA^{t}_{3}\bA^{t}_{2}\big),
\]
using $(\bA\bB)^{t}=\bB^{t}\bA^{t}$.

For the first equality, move $\bA^{t}_{2}$ to the front by cyclic
permutation and regroup:
\[
   \tr\big(\bA_{1}\bA^{t}_{3}\bA^{t}_{2}\big)
   =\tr\big(\bA^{t}_{2}\bA_{1}\bA^{t}_{3}\big)
   =\tr\big((\bA^{t}_{2}\bA_{1})\bA^{t}_{3}\big)
   =\bA^{t}_{2}\bA_{1}\cdot\bA_{3} .
\]

For the second, use $\tr\bX=\tr\bX^{t}$ on the same expression:
\[
\begin{aligned}
   \tr\big(\bA_{1}\bA^{t}_{3}\bA^{t}_{2}\big)
   &=\tr\big(\bA_{2}\bA_{3}\bA^{t}_{1}\big)
   =\tr\big(\bA_{3}\bA^{t}_{1}\bA_{2}\big)\\
   &=\tr\big((\bA_{3}\bA^{t}_{1})(\bA^{t}_{2})^{t}\big)
   =\bA_{3}\bA^{t}_{1}\cdot\bA^{t}_{2} ,
\end{aligned}
\]
the second step cyclic again. Hence
\[
   \boxed{\;
\begin{aligned}
   \bA_{1}\cdot\bA_{2}\bA_{3}
   &=\bA^{t}_{2}\bA_{1}\cdot\bA_{3}\\
   &=\bA_{3}\bA^{t}_{1}\cdot\bA^{t}_{2}
\end{aligned}\;}
\]

The rules say that a factor may be moved from one side of an inner
product to the other at the cost of a transpose, exactly as
$\bA\bu\cdot\bv=\bu\cdot\bA^{t}\bv$ does for vectors. They are used at
\eqr{11.22} and again in Problem 12 of Chapter 6.
\end{solution}

\begin{problem}[The isotropic Cauchy stress]
Derive \eqr{11.34}.
\end{problem}

\begin{solution}
Start from \eqr{11.33} and push it forward with \eqr{6.145},
$J\bT=\bF\bS\bF^{t}$:
\[
   J\bT=2\bF\big[\big(\hat W_{1}+I_{1}\hat W_{2}\big)\hI
   -\hat W_{2}\bC+I_{3}\hat W_{3}\bC^{-1}\big]\bF^{t},
\]
the scalar coefficients passing through $\bF$ untouched. Three products
are needed, and each is a one-line computation.

The identity term. $\bF\hI\bF^{t}=\bF\bF^{t}=\bB$, by the definition of
the left Cauchy--Green tensor at \eqr{2.83}.

The $\bC$ term. Substitute $\bC=\bF^{t}\bF$ and regroup:
\[
   \bF\bC\bF^{t}=\bF\big(\bF^{t}\bF\big)\bF^{t}
   =\big(\bF\bF^{t}\big)\big(\bF\bF^{t}\big)=\bB^{2}.
\]

The $\bC^{-1}$ term. Inverting $\bC=\bF^{t}\bF$ gives
$\bC^{-1}=\bF^{-1}\bF^{-t}$, so
\[
   \bF\bC^{-1}\bF^{t}=\bF\bF^{-1}\bF^{-t}\bF^{t}=\I\I=\I .
\]

Assembling,
\[
   \boxed{\;J\bT=2\big[\big(\hat W_{1}+I_{1}\hat W_{2}\big)\bB
   -\hat W_{2}\bB^{2}+I_{3}\hat W_{3}\I\big]\;}
\]

The pattern is worth noticing. Each referential tensor in \eqr{11.33} is
raised by one power of $\bB$: $\hI\mapsto\bB$, $\bC\mapsto\bB^{2}$,
$\bC^{-1}\mapsto\I$. That is the general rule $\bF\bC^{n}\bF^{t}
=\bB^{n+1}$, valid for every integer $n$, positive or negative, and it
is why the isotropic stress is a quadratic polynomial in $\bB$ exactly
when the energy gradient is a linear one in $\bC$ and $\bC^{-1}$.
\end{solution}

\begin{problem}[$\bC$ and $\bB$ have the same invariants]
Prove \eqr{11.35}.
\end{problem}

\begin{solution}
\textbf{By the spectral representations.} The right Cauchy--Green
tensor has the representation \eqr{3.38}$_{1}$,
\[
   \bC=\sum_{i}\lambda^{2}_{i}\,\bu_{i}\otimes\bu_{i},
   \qquad\text{and}\qquad
   \bB=\sum_{i}\lambda^{2}_{i}\,\bv_{i}\otimes\bv_{i},
\]
the $\bu_{i}$ being the referential principal axes, the
$\bv_{i}=\bR\bu_{i}$ the spatial ones, and the $\lambda_{i}$ the
principal stretches, common to both by
Theorem~\ref{thm:polar}. The principal invariants \eqr{11.31} are
symmetric functions of the eigenvalues,
\[
\begin{aligned}
   I_{1}&=\lambda^{2}_{1}+\lambda^{2}_{2}+\lambda^{2}_{3},\\
   I_{2}&=\lambda^{2}_{1}\lambda^{2}_{2}+\lambda^{2}_{2}\lambda^{2}_{3}
   +\lambda^{2}_{1}\lambda^{2}_{3},\\
   I_{3}&=\lambda^{2}_{1}\lambda^{2}_{2}\lambda^{2}_{3},
\end{aligned}
\]
and these expressions mention no eigenvector. Since $\bC$ and $\bB$
carry the same $\lambda^{2}_{i}$, they carry the same invariants:
\[
   \boxed{\;I_{k}(\bC)=I_{k}(\bB),\qquad k=1,2,3 .\;}
\]

\textbf{By similarity, in one line.} Substituting the polar
decomposition $\bF=\bR\bU$ into $\bB=\bF\bF^{t}$,
\[
   \bB=\bR\bU\bU^{t}\bR^{t}=\bR\bU^{2}\bR^{t}=\bR\bC\bR^{t},
\]
using $\bU^{t}=\bU$ and $\bC=\bU^{2}$ from \eqr{2.83}. So $\bB$ and
$\bC$ are related by a similarity transformation, and similar tensors
have the same characteristic polynomial:
\[
\begin{aligned}
   \det(\bB-\omega\I)&=\det\big(\bR(\bC-\omega\hI)\bR^{t}\big)\\
   &=\det\bR\,\det(\bC-\omega\hI)\,\det\bR^{t}
   =\det(\bC-\omega\hI),
\end{aligned}
\]
since $\det\bR\det\bR^{t}=\det(\bR\bR^{t})=1$. The invariants are the
coefficients of that polynomial, by \eqr{3.2}, so they agree. This
version does not need the spectral theorem, and it makes clear what is
really going on: $\bB$ is $\bC$ seen from the deformed configuration,
and a change of orthonormal frame cannot alter an invariant.
\end{solution}

\begin{problem}[The isotropic response as $O^{+}$ writes it]
Use \eqr{11.34}, or \eqr{11.36}, together with \eqr{11.49} to find
expressions for the stress and the heat flux in the frame of an
observer $O^{+}$.
\end{problem}

\begin{solution}
Chapter 7 supplies the transformations. Under a change of observer the
deformation gradient becomes $\bF^{+}=\bQ\bF$, with $\bQ(t)$ the
relative rotation, and the spatial temperature gradient becomes
$\grad^{+}\theta^{+}=\bQ\grad\theta$, while $\theta$ itself and the
reference configuration are untouched.

\textbf{What is invariant.} First,
\[
   \bB^{+}=\bF^{+}(\bF^{+})^{t}=\bQ\bF\bF^{t}\bQ^{t}=\bQ\bB\bQ^{t},
   \qquad
   \bC^{+}=(\bQ\bF)^{t}(\bQ\bF)=\bF^{t}\bF=\bC ,
\]
so $\bC$ does not move at all and $\bB$ moves by conjugation. Second,
$J^{+}=\det\bF^{+}=\det\bQ\det\bF=J$. Third, the invariants are
unchanged: $I_{k}(\bB^{+})=I_{k}(\bQ\bB\bQ^{t})=I_{k}(\bB)$ by the
similarity argument of Problem 3, and for the three of \eqr{11.50},
\[
\begin{aligned}
   I_{4}(\bB^{+},\grad^{+}\theta^{+})
   &=\ip{\bQ\grad\theta}{\bQ\bB\bQ^{t}\bQ\grad\theta}
   =\ip{\grad\theta}{\bB\grad\theta}=I_{4},
\end{aligned}
\]
using $\bQ^{t}\bQ=\I$ twice and $\ip{\bQ\ba}{\bQ\bb}=\ip\ba\bb$;
identically $I_{5}$ is unchanged, and $I_{6}=|\grad\theta|$ is unchanged
because $\bQ$ preserves length. So every one of the response functions
$\hat W_{k}$ and $\phi_{k}$ takes the same value for both observers.

\textbf{The stress.} Writing \eqr{11.34} in the frame of $O^{+}$ and
substituting,
\[
\begin{aligned}
   J\bT^{+}&=2\big[\big(\hat W_{1}+I_{1}\hat W_{2}\big)\bB^{+}
   -\hat W_{2}(\bB^{+})^{2}+I_{3}\hat W_{3}\I\big]\\
   &=2\big[\big(\hat W_{1}+I_{1}\hat W_{2}\big)\bQ\bB\bQ^{t}
   -\hat W_{2}\bQ\bB^{2}\bQ^{t}+I_{3}\hat W_{3}\bQ\bQ^{t}\big]
   =\bQ\big(J\bT\big)\bQ^{t},
\end{aligned}
\]
where $(\bB^{+})^{2}=\bQ\bB\bQ^{t}\bQ\bB\bQ^{t}=\bQ\bB^{2}\bQ^{t}$ and
$\I=\bQ\bQ^{t}$. Cancelling $J$,
\[
   \boxed{\;\bT^{+}=\bQ\bT\bQ^{t}\;}
\]
The incompressible form \eqr{11.36} gives the same conclusion, the
reaction transforming as $p^{+}=p$ because $\bQ(-p\I)\bQ^{t}=-p\I$.

\textbf{The heat flux.} By \eqr{11.49},
\[
\begin{aligned}
   \bq^{+}&=\big(\phi_{0}\I+\phi_{1}\bB^{+}+\phi_{2}(\bB^{+})^{2}\big)
   \grad^{+}\theta^{+}\\
   &=\bQ\big(\phi_{0}\I+\phi_{1}\bB+\phi_{2}\bB^{2}\big)\bQ^{t}
   \bQ\grad\theta
   =\bQ\,\bq .
\end{aligned}
\]
So
\[
   \boxed{\;\bq^{+}=\bQ\bq\;}
\]

\textbf{Reading the answer.} Nothing new has been derived. These are
precisely the transformation rules that Chapter 7 established for the
Cauchy stress and the heat flux from the balance laws, and the point of
the exercise is that the constitutive equations reproduce them
automatically. That is what frame invariance was imposed for: had
\eqr{11.34} delivered anything other than $\bQ\bT\bQ^{t}$, the two
observers would have disagreed about the material, and the model would
be inadmissible.
\end{solution}

\begin{problem}[Are the laminate constraints compatible with isotropy?]
Are the constraints considered in Problems 3 and 4 of Chapter 10
compatible with isotropy?
\end{problem}

\begin{solution}
The test is the one applied to incompressibility in
Remark~\ref{rem:ch11constraint}: a constraint is compatible with
isotropy when it is unchanged by $\bC\mapsto\bR^{t}\bC\bR$ for every
$\bR\in\Orth$. The answer in both cases is no, and for the same reason.

\textbf{The inextensible laminate.} The constraints of Problem 3 of
Chapter 10 are
\[
   \ip{\hbE_{1}}{\bC\hbE_{1}}=\ip{\hbE_{2}}{\bC\hbE_{2}}=1,
   \qquad \ip{\hbE_{1}}{\bC\hbE_{2}}=0 .
\]
Replace $\bC$ by $\bR^{t}\bC\bR$ and move $\bR$ across the inner
product:
\[
   \ip{\hbE_{A}}{\bR^{t}\bC\bR\hbE_{B}}
   =\ip{\bR\hbE_{A}}{\bC\,\bR\hbE_{B}} .
\]
The transformed constraints are therefore the original ones written
about the rotated directions $\bR\hbE_{1}$ and $\bR\hbE_{2}$. They agree
with the originals for every $\bC$ only if $\bR$ maps the plane spanned
by $\hbE_{1},\hbE_{2}$ onto itself, and most rotations do not. Take
$\bR$ to be the quarter turn about $\hbE_{1}$, which sends
$\hbE_{2}\mapsto\hbE_{3}$: the transformed constraint
$\ip{\hbE_{3}}{\bC\hbE_{3}}=1$ forbids transverse stretch, which the
original explicitly permits. So the constraint set is not invariant, and
the material is not isotropic.

\textbf{The area-preserving laminate.} The constraint of Problem 4 of
Chapter 10 says that the sheets, whose reference normal is $\hbE_{3}$,
do not change area. By the Piola--Nanson formula \eqr{2.54} the area
ratio is $|\bF^{*}\hbE_{3}|$, and
\[
   |\bF^{*}\hbE_{3}|^{2}=\ip{\hbE_{3}}{\bF^{*t}\bF^{*}\hbE_{3}}
   =J^{2}\ip{\hbE_{3}}{\bF^{-1}\bF^{-t}\hbE_{3}}
   =I_{3}\,\ip{\hbE_{3}}{\bC^{-1}\hbE_{3}},
\]
using $\bF^{*}=J\bF^{-t}$ and $J^{2}=\det\bC=I_{3}$. The constraint is
therefore
\[
   I_{3}\,\ip{\hbE_{3}}{\bC^{-1}\hbE_{3}}=1 .
\]
Under $\bC\mapsto\bR^{t}\bC\bR$ the factor $I_{3}$ is invariant, but the
second factor becomes $\ip{\bR\hbE_{3}}{\bC^{-1}\bR\hbE_{3}}$, so again
the constraint is carried to the same statement about a different
direction. Only rotations fixing the axis $\hbE_{3}$ leave it alone.
Not isotropic.

\textbf{Why incompressibility is the exception.} Both of these
constraints name a direction, $\hbE_{3}$ or the plane normal to it, and
a named direction is exactly what a rotation of the reference
configuration can move. Incompressibility names none: $\det\bC=1$
mentions only an invariant. That is the whole content of
Remark~\ref{rem:ch11constraint}, and Example~\ref{ex:ch10fibre} said the
same thing physically. A fibre or a laminate can be seen, so the
undeformed body is not the same in every direction.
\end{solution}

\begin{problem}[The invariants in simple shear]
Verify \eqr{11.38}.
\end{problem}

\begin{solution}
From \eqr{3.54} the left Cauchy--Green tensor of simple shear has the
matrix
\[
   (B_{ij})=\begin{pmatrix}1+\gamma^{2}&\gamma&0\\
   \gamma&1&0\\ 0&0&1\end{pmatrix}.
\]

\textbf{The first invariant} is the trace, by \eqr{11.31}$_{1}$:
\[
   I_{1}=B_{11}+B_{22}+B_{33}=(1+\gamma^{2})+1+1=3+\gamma^{2}.
\]

\textbf{The second} needs $\tr(\bB^{2})$. From \eqr{11.37}, whose
correction is the subject of Remark~\ref{rem:ch11mis37},
\[
   \tr(\bB^{2})=\big[(1+\gamma^{2})^{2}+\gamma^{2}\big]
   +(1+\gamma^{2})+1 .
\]
Expanding $(1+\gamma^{2})^{2}=1+2\gamma^{2}+\gamma^{4}$ and adding,
\[
   \tr(\bB^{2})=1+2\gamma^{2}+\gamma^{4}+\gamma^{2}+1+\gamma^{2}+1
   =3+4\gamma^{2}+\gamma^{4}.
\]
Then \eqr{11.31}$_{2}$ gives
\[
   I_{2}=\tfrac12\big[(3+\gamma^{2})^{2}
   -\big(3+4\gamma^{2}+\gamma^{4}\big)\big]
   =\tfrac12\big[9+6\gamma^{2}+\gamma^{4}
   -3-4\gamma^{2}-\gamma^{4}\big]=3+\gamma^{2}.
\]

\textbf{The third} is a determinant, and simple shear is isochoric:
\[
   I_{3}=\det\bB=\det(\bF\bF^{t})=(\det\bF)^{2}=1 .
\]
Hence
\[
   \boxed{\;I_{1}(\bB)=I_{2}(\bB)=3+\gamma^{2},\qquad I_{3}(\bB)=1 .\;}
\]

The coincidence $I_{1}=I_{2}$ is not an accident of the algebra. For an
isochoric deformation $I_{3}=1$, so the principal stretches satisfy
$\lambda_{1}\lambda_{2}\lambda_{3}=1$, and then
\[
   I_{2}=\lambda^{2}_{1}\lambda^{2}_{2}+\lambda^{2}_{2}\lambda^{2}_{3}
   +\lambda^{2}_{3}\lambda^{2}_{1}
   =\lambda^{-2}_{3}+\lambda^{-2}_{1}+\lambda^{-2}_{2},
\]
which is $I_{1}$ computed for the inverse deformation. Simple shear
happens to have the same stretches as its own inverse, up to a
relabelling, which is why the two invariants agree here.
\end{solution}

\begin{problem}[The universal relation]
Derive \eqr{11.42}.
\end{problem}

\begin{solution}
Take the $11$ and $22$ components of \eqr{11.39}. The first matrix
contributes $f_{0}$ to each, so it cancels in the difference. The other
two contribute
\[
\begin{aligned}
   T_{11}&=f_{0}+f_{1}(1+\gamma^{2})
   +f_{2}\big[(1+\gamma^{2})^{2}+\gamma^{2}\big],\\
   T_{22}&=f_{0}+f_{1}\cdot1+f_{2}(1+\gamma^{2}).
\end{aligned}
\]
Subtracting,
\[
\begin{aligned}
   T_{11}-T_{22}
   &=f_{1}\big[(1+\gamma^{2})-1\big]
    +f_{2}\big[(1+\gamma^{2})^{2}+\gamma^{2}-(1+\gamma^{2})\big]\\
   &=f_{1}\gamma^{2}
    +f_{2}\big[(1+\gamma^{2})\big((1+\gamma^{2})-1\big)+\gamma^{2}\big]\\
   &=f_{1}\gamma^{2}+f_{2}\big[(1+\gamma^{2})\gamma^{2}+\gamma^{2}\big]
    =\gamma^{2}\big[f_{1}+(2+\gamma^{2})f_{2}\big],
\end{aligned}
\]
the last step taking $\gamma^{2}$ out and simplifying
$(1+\gamma^{2})+1=2+\gamma^{2}$. The bracket is $\mu(\gamma^{2})$ by
\eqr{11.40}, so $T_{11}-T_{22}=\gamma^{2}\mu$.

On the other hand $T_{12}=\gamma\mu$, again by \eqr{11.40}, so
$\gamma T_{12}=\gamma^{2}\mu$ as well. Hence
\[
   \boxed{\;T_{11}-T_{22}=\gamma\,T_{12}\;}
\]

Two remarks on the mechanism. The constitutive functions entered only
through the combination $\mu$, and they entered the two sides in the
same way, which is why they cancel; that cancellation is the whole
content of the word \emph{universal}. And the computation used the
corrected $11$ entry of $\bB^{2}$ throughout: with the printed value the
$f_{2}$ bracket collapses to $\gamma^{2}$ instead of
$\gamma^{2}(2+\gamma^{2})$, and the relation fails, which is the
evidence given in Remark~\ref{rem:ch11mis37}.
\end{solution}

\begin{problem}[The neo-Hookean material, and a bar in tension]
Consider the strain energy $W^{*}=\tfrac12 G(I_{1}-3)$ with $G$ a
positive constant. This model, the \emph{neo-Hookean} material,
describes incompressible rubber well while the principal stretches stay
in the approximate range $\tfrac12<\lambda_{i}<2$.
\textup{(a)} Show that $\bT=-p\I+G\bB$, and that equilibrium without
body force requires $\grad p=G\dvg\bB$.
\textup{(b)} Use the power conjugacy of $\bP$ and $\bF$ to deduce that
$G$ is the ratio of the shear component of the Piola traction to the
amount of shear in a simple shear.
\textup{(c)} For the deformation of Problem 6 of Chapter 4, find $\bC$
and $\bB$.
\textup{(d)} Impose $\bT\bn=\bze$ for every unit $\bn$ with
$\ip\bn{\be_{1}}=0$, simulating the traction-free lateral surface of a
prismatic bar with axis $\be_{1}$, and deduce $p=G/\lambda$.
\textup{(e)} Show that $\bT=T(\lambda)\be_{1}\otimes\be_{1}$ with
$T(\lambda)=G(\lambda^{2}-\lambda^{-1})$.
\textup{(f)} Young's modulus is $E=T'(1)$. Find $E$ in terms of $G$.
\textup{(g)} Show that $\bP=P(\lambda)\be_{1}\otimes\hbE_{1}$ and
$\bS=S(\lambda)\hbE_{1}\otimes\hbE_{1}$, determine $P$ and $S$, and
relate $P'(1)$ and $S'(1)$ to $E$.
\end{problem}

\begin{solution}
\textbf{(a)} With $W^{*}=\tfrac12G(I_{1}-3)$ the derivatives are
$W^{*}_{1}=\tfrac12G$ and $W^{*}_{2}=0$, so \eqr{11.36} gives
\[
   \bT=-p\I+2\big[\big(\tfrac12G+I_{1}\cdot0\big)\bB
   -0\cdot\bB^{2}\big]=\boxed{\;-p\I+G\bB\;}
\]
Equilibrium without body force is $\dvg\bT=\bze$ by \eqr{6.109} with
$\dot\bv=\bze$. Since $\dvg(p\I)=\grad p$, by Problem 24 of Chapter 1,
and $G$ is constant,
\[
   -\grad p+G\dvg\bB=\bze,\qquad\text{that is}\qquad
   \boxed{\;\grad p=G\,\dvg\bB\;}
\]
The equation determines $p$ up to a constant once the deformation is
known, and it is solvable only if $\dvg\bB$ is a gradient. That is a
genuine restriction: not every isochoric deformation of a neo-Hookean
body can be maintained without body force.

\textbf{(b)} By \eqr{6.156} the stress power per unit reference volume
is $\bP\cdot\dot\bF$, and for a hyperelastic material at fixed
temperature that equals $\dot W^{*}$. In simple shear
$\bF=\shift+\gamma\be_{1}\otimes\hbE_{2}$ with $I_{1}=3+\gamma^{2}$ by
\eqr{11.38}, so
\[
   W^{*}=\tfrac12G\big(3+\gamma^{2}-3\big)=\tfrac12G\gamma^{2},
   \qquad \dot W^{*}=G\gamma\dot\gamma .
\]
On the other side, $\dot\bF=\dot\gamma\,\be_{1}\otimes\hbE_{2}$, and by
Definition~\ref{def:tip},
\[
   \bP\cdot\dot\bF=\dot\gamma\,\bP\cdot\be_{1}\otimes\hbE_{2}
   =\dot\gamma\,\ip{\be_{1}}{\bP\hbE_{2}}=\dot\gamma\,P_{12}.
\]
The reactive part contributes nothing, since its share of the power is
$-p\bF^{*}\cdot\dot\bF=-p\dot J=0$ by Problem 2 of Chapter 4 and
$J\equiv1$. Equating and cancelling the arbitrary $\dot\gamma$,
\[
   \boxed{\;P_{12}=G\gamma,\qquad G=P_{12}/\gamma\;}
\]
and $P_{12}=\ip{\be_{1}}{\bP\hbE_{2}}$ is the $\be_{1}$ component of the
Piola traction on the material plane whose reference normal is
$\hbE_{2}$, which is the shear component. So $G$ is the shear stress per
unit shear, the shear modulus, and here the identification is exact
rather than a small-strain approximation.

\textbf{(c)} Problem 6 of Chapter 4 has
\[
   \bF=\lambda\,\be_{1}\otimes\hbE_{1}
   +\lambda^{-1/2}\big(\be_{2}\otimes\hbE_{2}+\be_{3}\otimes\hbE_{3}\big),
\]
an axial stretch $\lambda$ with equal lateral contractions, isochoric
because $\lambda\cdot\lambda^{-1/2}\cdot\lambda^{-1/2}=1$. Multiplying
out with $(\ba\otimes\bb)(\bc\otimes\bd)=(\ip\bb\bc)\ba\otimes\bd$, only
the diagonal terms surviving,
\[
   \boxed{\;
\begin{aligned}
   \bC&=\bF^{t}\bF=\lambda^{2}\hbE_{1}\otimes\hbE_{1}
   +\lambda^{-1}\big(\hbE_{2}\otimes\hbE_{2}
   +\hbE_{3}\otimes\hbE_{3}\big),\\
   \bB&=\bF\bF^{t}=\lambda^{2}\be_{1}\otimes\be_{1}
   +\lambda^{-1}\big(\be_{2}\otimes\be_{2}+\be_{3}\otimes\be_{3}\big).
\end{aligned}\;}
\]
They carry the same numbers on different bases, as Problem 3 requires.

\textbf{(d)} Let $\bn$ be a unit vector with $\ip\bn{\be_{1}}=0$, so
that $\bn$ lies in the plane spanned by $\be_{2}$ and $\be_{3}$. On that
plane $\bB$ acts as $\lambda^{-1}$ times the identity, by (c), so
$\bB\bn=\lambda^{-1}\bn$ and
\[
   \bT\bn=-p\bn+G\lambda^{-1}\bn
   =\big(-p+G\lambda^{-1}\big)\bn .
\]
Setting this to zero, and $\bn\neq\bze$,
\[
   \boxed{\;p=G/\lambda\;}
\]
The pressure is uniform, so $\grad p=\bze$; and $\bB$ is uniform, so
$\dvg\bB=\bze$; hence (a) is satisfied and the deformation is
equilibrated without body force.

\textbf{(e)} Substituting $p=G/\lambda$ into (a) and using (c),
\[
\begin{aligned}
   \bT&=-\frac{G}{\lambda}\I+G\Big[\lambda^{2}\be_{1}\otimes\be_{1}
   +\lambda^{-1}\big(\be_{2}\otimes\be_{2}
   +\be_{3}\otimes\be_{3}\big)\Big]\\
   &=G\Big(\lambda^{2}-\frac{1}{\lambda}\Big)\be_{1}\otimes\be_{1}
   +G\Big(\frac{1}{\lambda}-\frac{1}{\lambda}\Big)
   \big(\be_{2}\otimes\be_{2}+\be_{3}\otimes\be_{3}\big),
\end{aligned}
\]
the lateral coefficients cancelling exactly, which is what (d) arranged.
Hence
\[
   \boxed{\;\bT=T(\lambda)\,\be_{1}\otimes\be_{1},\qquad
   T(\lambda)=G\big(\lambda^{2}-\lambda^{-1}\big).\;}
\]
The bar carries a uniaxial stress, and $T(1)=0$: the undeformed state is
unstressed, as it must be for a model with no residual stress.

\textbf{(f)} Differentiating,
$T'(\lambda)=G(2\lambda+\lambda^{-2})$, so
\[
   \boxed{\;E=T'(1)=G(2+1)=3G\;}
\]
This is the incompressible limit of the classical relation
$E=2G(1+\nu)$ at $\nu=\tfrac12$, and Chapter 12 will recover it from the
linear theory.

\textbf{(g)} By \eqr{6.140}, $\bP=\bT\bF^{*}=J\bT\bF^{-t}$ with $J=1$.
Problem 6(b) of Chapter 4 gives
$\bF^{-1}=\lambda^{-1}\hbE_{1}\otimes\be_{1}
+\lambda^{1/2}(\hbE_{2}\otimes\be_{2}+\hbE_{3}\otimes\be_{3})$, so
\[
   \bF^{-t}=\lambda^{-1}\be_{1}\otimes\hbE_{1}
   +\lambda^{1/2}\big(\be_{2}\otimes\hbE_{2}
   +\be_{3}\otimes\hbE_{3}\big).
\]
Multiplying, only the $\be_{1}$ term of $\bT$ surviving,
\[
   \bP=T(\lambda)\big(\be_{1}\otimes\be_{1}\big)\bF^{-t}
   =T(\lambda)\lambda^{-1}\,\be_{1}\otimes\hbE_{1},
\]
and by \eqr{6.144}, $\bS=\bF^{-1}\bP=P(\lambda)\lambda^{-1}
\hbE_{1}\otimes\hbE_{1}$. So
\[
   \boxed{\;P(\lambda)=\frac{T(\lambda)}{\lambda}
   =G\big(\lambda-\lambda^{-2}\big),\qquad
   S(\lambda)=\frac{P(\lambda)}{\lambda}
   =G\big(1-\lambda^{-3}\big).\;}
\]
Differentiating, $P'(\lambda)=G(1+2\lambda^{-3})$ and
$S'(\lambda)=3G\lambda^{-4}$, so
\[
   P'(1)=3G=E,\qquad S'(1)=3G=E .
\]
All three stress measures have the same slope at the undeformed state.
That is the general reason why the linear theory of Chapter 12 need not
say which stress it means: the three differ at second order in the
strain, and the linear theory keeps only the first.
\end{solution}

\begin{problem}[Anti-plane shear of a neo-Hookean tube]
A circular cylindrical bar of neo-Hookean material, with axis $\hbE_{3}$
in $\kappa$ and inner and outer radii $A$ and $B$, undergoes the static
anti-plane shear of Problem 5 of Chapter 3. It is in equilibrium under
negligible body force, and $w(A)=W$, $w(B)=0$ with $W$ a specified
constant. Obtain $w(R)$.
\end{problem}

\begin{solution}
\textbf{The kinematics.} The deformation is $r=R$, $\theta=\Theta$,
$z=Z+w(R)$: cylinders slide along the axis without changing radius.
Identifying $\hbE_{R}$ with $\ber$ and $\hbE_{\Theta}$ with $\bet$
through the shifter, since $r=R$ and $\theta=\Theta$,
\[
   \bF=\shift+w'(r)\,\be_{3}\otimes\hbE_{R},\qquad \det\bF=1 ,
\]
the determinant being that of a matrix with unit diagonal and a single
entry below it. Then
\[
   \bB=\bF\bF^{t}=\I+w'\big(\ber\otimes\be_{3}+\be_{3}\otimes\ber\big)
   +w'^{2}\,\be_{3}\otimes\be_{3},
\]
so the non-zero components are $B_{rr}=B_{\theta\theta}=1$,
$B_{33}=1+w'^{2}$ and $B_{r3}=B_{3r}=w'$.

\textbf{The stress.} By Problem 8(a),
\[
   T_{rr}=T_{\theta\theta}=-p+G,\qquad T_{33}=-p+G(1+w'^{2}),
   \qquad T_{r3}=G\,w'(r),
\]
all functions of $r$ alone, the remaining components zero.

\textbf{Equilibrium.} With nothing depending on $\theta$ or $z$, the
cylindrical components of $\dvg\bT$ reduce to
\[
   (\dvg\bT)_{r}=\frac{\partial T_{rr}}{\partial r}
   +\frac{T_{rr}-T_{\theta\theta}}{r},
   \qquad
   (\dvg\bT)_{3}=\frac{1}{r}\frac{\partial}{\partial r}\big(rT_{3r}\big),
\]
and $(\dvg\bT)_{\theta}=0$ identically. Since $T_{rr}=T_{\theta\theta}$,
the radial equation is $-p'(r)=0$, so the reactive pressure is uniform.
The axial equation is the one that determines the motion:
\[
   \frac{1}{r}\big(G\,r\,w'\big)'=0
   \qquad\Longrightarrow\qquad r\,w'(r)=c
\]
for a constant $c$, whence $w'=c/r$ and
\[
   w(r)=c\ln r+d .
\]

\textbf{The boundary conditions.} Imposing $w(A)=W$ and $w(B)=0$,
\[
   c\ln A+d=W,\qquad c\ln B+d=0 ,
\]
so $d=-c\ln B$ and $c\ln(A/B)=W$. Therefore
\[
   \boxed{\;w(R)=W\,\frac{\ln(R/B)}{\ln(A/B)}\;}
\]
which returns $W$ at $R=A$ and $0$ at $R=B$, as required.

\textbf{Reading the answer.} The profile is logarithmic, and it contains
no material constant: $G$ cancelled when the axial equation was divided
through. So every neo-Hookean tube, stiff or soft, shears into the same
shape, and $G$ affects only the force needed. The axial traction on the
inner surface is $T_{3r}(A)=Gc/A=GW/[A\ln(A/B)]$, and the total axial
force there, over a length $l$, is $2\pi A l\,T_{3r}(A)
=2\pi lGW/\ln(A/B)$, independent of $A$; the same magnitude with
opposite sign acts on the outer surface, as it must for equilibrium of
the tube.
\end{solution}

\begin{problem}[Residual stress in an isotropic solid]
An elastic solid is isotropic at all material points relative to the
adopted reference configuration, and is stressed in its undeformed
state, so that $\bT\neq\Ob$ when $\bF=\shift$. Suppose this residually
stressed undeformed configuration is in equilibrium without body force,
and that the traction vanishes on a portion of the boundary. Show that
the residual stress vanishes identically.
\end{problem}

\begin{solution}
\textbf{Step 1: isotropy forces the residual stress to be a pressure.}
At $\bF=\shift$ the left Cauchy--Green tensor is
$\bB=\shift\shift^{t}=\I$, so
\[
   I_{1}=\tr\I=3,\qquad I_{2}=\tfrac12\big[9-\tr\I\big]=3,
   \qquad I_{3}=\det\I=1 ,
\]
and $\bB^{2}=\I$. Substituting into \eqr{11.34}, with $J=1$,
\[
   \bT=2\big[\big(\hat W_{1}+3\hat W_{2}\big)\I
   -\hat W_{2}\I+\hat W_{3}\I\big]
   =2\big[\hat W_{1}+2\hat W_{2}+\hat W_{3}\big]\I ,
\]
every derivative evaluated at $(\theta,3,3,1)$. So
\[
   \bT=-p_{0}(\bX)\,\I,\qquad
   p_{0}=-2\big[\hat W_{1}+2\hat W_{2}+\hat W_{3}\big]\big|_{(3,3,1)} .
\]
This is the whole force of isotropy: whatever the material, a residual
stress compatible with isotropy can only be hydrostatic. It may still
vary from point to point, since $\hat W$ may depend on $\bX$.

\textbf{Step 2: equilibrium forces it to be uniform.} With no body force
and no acceleration, \eqr{6.109} reads $\dvg\bT=\bze$, and by Problem 24
of Chapter 1,
\[
   \dvg(-p_{0}\I)=-\grad p_{0}=\bze ,
\]
so $p_{0}$ is constant throughout the body, the body being connected.

\textbf{Step 3: the free boundary forces it to be zero.} Let $\bn$ be
the unit normal at a point of the traction-free portion. There
\[
   \bze=\bT\bn=-p_{0}\bn ,
\]
and $|\bn|=1$, so $p_{0}=0$. Being constant, $p_{0}$ vanishes
everywhere, and therefore
\[
   \boxed{\;\bT=\Ob\ \text{ throughout the body.}\;}
\]

\textbf{What each hypothesis did.} Drop isotropy and Step 1 fails: an
anisotropic solid can carry a residual stress with a non-zero deviator,
which is what a shrink-fit ring or a cold-drawn wire does, and such
stresses are the reason residual stress is a subject at all. Drop
equilibrium and Step 2 fails. Drop the traction-free patch and Step 3
fails, leaving an arbitrary uniform pressure, which is exactly the
indeterminacy met in Problem 13(e) below. So the result is sharp: an
isotropic body cannot hold residual stress unless something holds it
from outside.
\end{solution}

\begin{problem}[An azimuthal anti-plane shear, and what it is]
Consider the deformation $\bx=\shift\bX+w(\Theta)\be_{3}$, with $\Theta$
the azimuthal angle of a cylindrical polar system $\{R,\Theta,Z\}$ in
the reference configuration.
\textup{(a)} Using $\bx=r\ber(\theta)+z\be_{3}$ with a suitable choice
of basis vectors, show that one may take $R=r$, $\Theta=\theta$ and
$Z+w=z$ without loss of generality.
\textup{(b)} Show that
$\bF=\shift+R^{-1}w'(\Theta)\,\be_{3}\otimes\hbE_{\Theta}$, obtain
$\bB$, and show $\det\bF=1$.
\textup{(c)} If a neo-Hookean material occupies an annular tube of
radii $A$ and $B$, is in equilibrium with zero body force, and has
traction-free inner and outer cylindrical surfaces, show that
$w(\Theta)=C\Theta+D$. Interpret the deformation.
\end{problem}

\begin{solution}
\textbf{(a)} The shifter carries $\hbE_{R}(\Theta)\mapsto\ber(\Theta)$,
$\hbE_{\Theta}\mapsto\bet$ and $\hbE_{3}\mapsto\be_{3}$, so with
$\bX=R\hbE_{R}(\Theta)+Z\hbE_{3}$,
\[
   \bx=\shift\bX+w(\Theta)\be_{3}
   =R\,\ber(\Theta)+\big(Z+w(\Theta)\big)\be_{3}.
\]
Comparing with $\bx=r\ber(\theta)+z\be_{3}$, and using that
$\{\ber,\bet,\be_{3}\}$ is a basis at each point, the radial components
give $r=R$ and the polar angle of $\ber$ gives $\theta=\Theta$, while
the axial components give $z=Z+w$.

\textbf{(b)} Write $\bx=\shift\bX+w(\Theta)\be_{3}$ and take the
referential gradient term by term. The first term contributes $\shift$,
since $\Grad(\shift\bX)=\shift$. For the second, $\be_{3}$ is constant
and $w$ is a scalar field, so $\Grad(w\be_{3})=\be_{3}\otimes\Grad w$ by
Problem 24(f) of Chapter 1, and the cylindrical gradient of a function
of $\Theta$ alone is $\Grad w=R^{-1}w'(\Theta)\hbE_{\Theta}$. Hence
\[
   \boxed{\;\bF=\shift+\alpha\,\be_{3}\otimes\hbE_{\Theta},
   \qquad \alpha=\frac{w'(\Theta)}{R} .\;}
\]
On the ordered bases $\{\hbE_{R},\hbE_{\Theta},\hbE_{3}\}$ and
$\{\ber,\bet,\be_{3}\}$ its matrix is
$\left(\begin{smallmatrix}1&0&0\\0&1&0\\0&\alpha&1\end{smallmatrix}
\right)$, lower triangular with unit diagonal, so $\det\bF=1$ and the
deformation is isochoric. For the left Cauchy--Green tensor,
\[
\begin{aligned}
   \bB&=\bF\bF^{t}
   =\big(\shift+\alpha\be_{3}\otimes\hbE_{\Theta}\big)
    \big(\shift^{t}+\alpha\hbE_{\Theta}\otimes\be_{3}\big)\\
   &=\shift\shift^{t}
    +\alpha\big(\shift\hbE_{\Theta}\big)\otimes\be_{3}
    +\alpha\,\be_{3}\otimes\big(\shift\hbE_{\Theta}\big)
    +\alpha^{2}\big(\ip{\hbE_{\Theta}}{\hbE_{\Theta}}\big)
     \be_{3}\otimes\be_{3}\\
   &=\I+\alpha\big(\bet\otimes\be_{3}+\be_{3}\otimes\bet\big)
    +\alpha^{2}\,\be_{3}\otimes\be_{3},
\end{aligned}
\]
using $\shift\shift^{t}=\I$ and $\shift\hbE_{\Theta}=\bet$.

\textbf{(c)} By Problem 8(a) the stress is $\bT=-p\I+G\bB$, so
\[
   T_{rr}=T_{\theta\theta}=-p+G,\qquad
   T_{33}=-p+G(1+\alpha^{2}),\qquad
   T_{\theta3}=G\alpha ,
\]
the remaining components zero.

\emph{The lateral surfaces.} On $r=A$ and $r=B$ the normal is $\pm\ber$,
and $\bT\ber=T_{rr}\ber$ since $T_{r\theta}=T_{r3}=0$. Traction-free
therefore means $p=G$ on both.

\emph{The radial equation.} With $T_{r\theta}=0$ and
$T_{rr}=T_{\theta\theta}$,
\[
   (\dvg\bT)_{r}=\frac{\partial T_{rr}}{\partial r}
   +\frac{T_{rr}-T_{\theta\theta}}{r}
   =-\frac{\partial p}{\partial r}=0 ,
\]
so $p$ does not depend on $r$. Combined with $p=G$ at $r=A$, this gives
$p\equiv G$ throughout, and then $T_{rr}=T_{\theta\theta}=0$ everywhere.

\emph{The axial equation.} With $T_{3r}=0$ and $T_{33}$ independent of
$z$,
\[
   (\dvg\bT)_{3}=\frac{1}{r}\frac{\partial T_{3\theta}}{\partial\theta}
   =\frac{G}{r}\frac{\partial\alpha}{\partial\theta}=0 ,
\]
so $\alpha=w'(\theta)/r$ must not depend on $\theta$. Hence
$w'(\Theta)$ is a constant $C$ and
\[
   \boxed{\;w(\Theta)=C\Theta+D\;}
\]
The azimuthal equation is satisfied identically, every relevant
component being independent of $\theta$.

\textbf{Interpretation.} The deformation is $z=Z+C\Theta$: material at
azimuth $\Theta$ is lifted along the axis in proportion to $\Theta$.
Going once around the tube, $\Theta\mapsto\Theta+2\pi$ returns to the
same material line but at an axial offset $2\pi C$, so $w$ is not single
valued on the tube unless $C=0$. The deformation cannot be produced from
the intact tube. It is produced by cutting the tube along a half plane
$\Theta=\text{const}$, sliding one face axially by $2\pi C$, and welding
the faces together again. That is a \emph{screw dislocation} of Burgers
vector $2\pi C\be_{3}$ along the axis, and \eqr{11.42} has nothing to do
with it: the stress computed above has $T_{rr}=T_{\theta\theta}=0$ and
the only non-zero component is the shear $T_{\theta3}=GC/r$, which is
singular on the axis and is why real dislocation cores need a separate
model. Remark~\ref{rem:P9multi} of Chapter 4 met the same
multivaluedness in a velocity field.
\end{solution}

\begin{problem}[Torsion of a neo-Hookean bar, and the Poynting effect]
A solid circular cylindrical bar of neo-Hookean material is subjected to
the torsion of Section~\ref{sec:torsion}. It is in equilibrium without
body force and the traction on its cylindrical surface vanishes. Obtain
the stress field and use it to compute the net force and torque on a
cross-section. Show that the torque and twist are related as classical
strength of materials predicts, that the axial force is non-zero
contrary to that theory, and that it is proportional to the square of
the twist.
\end{problem}

\begin{solution}
\textbf{The kinematics.} Torsion carries $\bX=R\hbE_{R}(\Theta)+Z
\hbE_{3}$ to $\bx=R\ber(\Theta+\tau Z)+Z\be_{3}$, with $\tau$ the twist
per unit length. Differentiating along the three referential
directions,
\[
   \frac{\partial\bx}{\partial R}=\ber,\qquad
   \frac{1}{R}\frac{\partial\bx}{\partial\Theta}=\bet,\qquad
   \frac{\partial\bx}{\partial Z}=\tau R\,\bet+\be_{3},
\]
so
\[
   \bF=\shift+\gamma\,\bet\otimes\hbE_{3},\qquad \gamma=\tau R=\tau r ,
\]
and $\det\bF=1$. The local shear is $\gamma$, growing linearly with
radius, which is the sense in which torsion is locally a simple shear.
Squaring as in Problem 11(b),
\[
   \bB=\I+\gamma\big(\bet\otimes\be_{3}+\be_{3}\otimes\bet\big)
   +\gamma^{2}\,\bet\otimes\bet .
\]

\textbf{The stress.} By Problem 8(a), $\bT=-p\I+G\bB$, so
\[
   T_{rr}=-p+G,\qquad T_{\theta\theta}=-p+G\big(1+\gamma^{2}\big),
   \qquad T_{zz}=-p+G,\qquad T_{\theta z}=G\gamma ,
\]
the rest zero. Note $T_{rr}=T_{zz}$ and
$T_{rr}-T_{\theta\theta}=-G\gamma^{2}=-G\tau^{2}r^{2}$.

\textbf{Equilibrium.} The azimuthal and axial components of $\dvg\bT$
vanish identically, nothing depending on $\theta$ or $z$ and $T_{zr}=0$.
The radial component gives
\[
   \frac{\partial T_{rr}}{\partial r}
   +\frac{T_{rr}-T_{\theta\theta}}{r}=0
   \qquad\Longrightarrow\qquad
   -p'(r)-G\tau^{2}r=0 ,
\]
so $p'(r)=-G\tau^{2}r$ and $p(r)=-\tfrac12G\tau^{2}r^{2}+c$.

\textbf{The free lateral surface.} At $r=a$ the normal is $\ber$ and
$\bT\ber=T_{rr}\ber$, so $T_{rr}(a)=0$, that is $p(a)=G$. Hence
$c=G+\tfrac12G\tau^{2}a^{2}$ and
\[
   \boxed{\;p(r)=G+\tfrac12G\tau^{2}\big(a^{2}-r^{2}\big)\;}
\]
Substituting back, the stress field is
\[
   \boxed{\;
\begin{aligned}
   T_{rr}&=T_{zz}=-\tfrac12G\tau^{2}\big(a^{2}-r^{2}\big),\\
   T_{\theta\theta}&=G\tau^{2}\big(\tfrac32r^{2}-\tfrac12a^{2}\big),
   \qquad T_{\theta z}=G\tau r .
\end{aligned}\;}
\]

\textbf{The torque.} The traction on a cross-section, whose normal is
$\be_{3}$, is $\bT\be_{3}=T_{\theta z}\bet+T_{zz}\be_{3}$. Its moment
about the axis is
\[
   M=\int_{\text{section}}\big[(r\ber)\times\bT\be_{3}\big]
   \cdot\be_{3}\,\dl a
   =\int^{a}_{0}\!\!\big(r\cdot G\tau r\big)\,2\pi r\,\dl r
   =2\pi G\tau\,\frac{a^{4}}{4},
\]
using $\ber\times\bet=\be_{3}$ and $\dl a=2\pi r\,\dl r$. Hence
\[
   \boxed{\;M=\tfrac12\pi G\,\tau\,a^{4}=G\,J_{p}\,\tau,
   \qquad J_{p}=\tfrac12\pi a^{4} ,\;}
\]
$J_{p}$ being the polar second moment of the section. This is exactly
the classical formula, with no nonlinear correction whatever: the
torque is linear in the twist however large the twist becomes. The
reason is visible in the computation: $T_{\theta z}=G\gamma$ exactly,
the neo-Hookean shear response being linear in the amount of shear.

\textbf{The axial force.} On the same section,
\[
   N=\int^{a}_{0}T_{zz}\,2\pi r\,\dl r
   =-\pi G\tau^{2}\int^{a}_{0}\big(a^{2}-r^{2}\big)r\,\dl r
   =-\pi G\tau^{2}\Big[\frac{a^{4}}{2}-\frac{a^{4}}{4}\Big],
\]
so
\[
   \boxed{\;N=-\tfrac14\pi G\,\tau^{2}a^{4}\;}
\]
It is not zero, it is compressive, and it is quadratic in $\tau$.
Classical strength of materials, being linear, predicts $N=0$.

\textbf{What it means.} To hold the bar at fixed length while twisting
it, one must push on the ends with the force $|N|$. Release the ends and
the bar lengthens instead. This is the \emph{Poynting effect}, and it is
the same phenomenon as the normal stress effect of
Remark~\ref{rem:ch11universal}: torsion is locally a simple shear of
amount $\gamma=\tau r$, and \eqr{11.42} demands normal stresses of order
$\gamma^{2}$ to accompany it. Integrating those over the section is
exactly the calculation just done. The ratio
\[
   \frac{|N|}{M/a}=\frac{\tfrac14\pi G\tau^{2}a^{4}}
   {\tfrac12\pi G\tau a^{3}}=\tfrac12\tau a
\]
is half the surface shear, so the effect becomes measurable when the
twist per unit length times the radius is not small.
\end{solution}

\bookfig{1}\begin{figure}[htbp]\centering
\begin{tikzpicture}[scale=1.0,
   ar/.style={-{Stealth[length=2.2mm]},thick},
   dim/.style={{Stealth[length=1.7mm]}-{Stealth[length=1.7mm]},
               semithick,gray!75}]
% ---------------- reference: a sector of angular extent gamma --------
\begin{scope}[xshift=0cm]
  \fill[gray!10] (0,0) -- (-30:1.85) arc (-30:240:1.85) -- cycle;
  \draw[thick] (0,0) -- (-30:1.85) arc (-30:240:1.85) -- cycle;
  % the wedge that was removed
  \draw[guide] (240:1.85) arc (240:330:1.85);
  \fill (0,0) circle (1.2pt);
  % the angular extent of the material
  \draw[gray!75,thin] (-30:0.55) arc (-30:240:0.55);
  \node[font=\small] at (150:0.80) {$\gamma$};
  % the reference radius
  \draw[dim] (0.05,0.05) -- (45:1.80);
  \node[font=\small,fill=white,inner sep=1.5pt] at (45:1.22) {$A$};
  \node[font=\footnotesize,text=gray!80,anchor=north] at (0,-2.22)
     {$\kappa$: a wedge is missing};
\end{scope}
% ---------------- deformed: the full disc ----------------------------
\begin{scope}[xshift=5.6cm]
  \fill[gray!10] (0,0) circle (1.52);
  \draw[thick] (0,0) circle (1.52);
  \draw[thick,gray!55] (0,0) -- (270:1.52);
  \fill (0,0) circle (1.2pt);
  \draw[dim] (0.05,0.05) -- (45:1.47);
  \node[font=\small,fill=white,inner sep=1.5pt] at (45:0.92) {$a$};
  \node[font=\footnotesize,text=gray!80,anchor=east] at (-0.14,-1.02)
     {seam};
  \node[font=\footnotesize,text=gray!80,anchor=north] at (0,-2.22)
     {$\kappa_{t}$: the gap is closed};
\end{scope}
\end{tikzpicture}
\caption{Problem 13. In the reference configuration the material of a
circular cylinder occupies an angular extent $\gamma<2\pi$, a wedge
having been removed; the missing wedge is dashed. The cross-section is
then deformed so that the two cut faces meet, the material spanning the
full $2\pi$. The map is $r=f(R)$, $\theta=g(\Theta)$, $z=Z$, and
isochoricity together with regularity on the axis fixes both functions.
The deformed radius comes out proportional to $\sqrt\gamma$, so closing
a larger gap makes a thinner bar.}
\label{fig:ch11wedge}
\end{figure}

\begin{problem}[Closing a wedge: the cylinder with a sector removed]
A circular cylinder of radius $A$ has a wedge removed, so that the
angular extent of the material in $\kappa$ is $\gamma$
(Figure~\ref{fig:ch11wedge}). The cross-section is deformed so that the
gap closes, by
\[
\begin{aligned}
   \bX&=R\hbE_{R}(\Theta)+Z\hbE_{3},\qquad
   \bx=r\ber(\theta)+z\be_{3},\\
   z&=Z,\qquad r=f(R),\qquad \theta=g(\Theta).
\end{aligned}
\]
\textup{(a)} Find $\bF$ for known $f$ and $g$.
\textup{(b)} Assume the deformation isochoric and find $f$ and $g$ from
the boundary conditions; show the deformed radius $a$ is proportional to
$\sqrt\gamma$.
\textup{(c)} Find $\bR$, $\bU$, $\bV$.
\textup{(d)} Obtain the principal stretches in terms of $\gamma$.
\textup{(e)} Solve the equilibrium problem without body force for a
neo-Hookean material. Find the reactive pressure and show it is singular
on the centreline. Evaluate the traction at the outer boundary required
to maintain the deformation and show it is an arbitrary uniform
pressure. What contact traction acts on the planar surfaces where the
two faces of the wedge meet, if there is no traction at the outer
boundary?
\end{problem}

\begin{solution}
\textbf{(a)} Differentiating $\bx$ along the three referential
directions, and remembering that $\ber$ depends on $\theta=g(\Theta)$
so that $\partial\ber/\partial\Theta=g'(\Theta)\bet$,
\[
   \frac{\partial\bx}{\partial R}=f'(R)\ber,\qquad
   \frac{1}{R}\frac{\partial\bx}{\partial\Theta}
   =\frac{f(R)g'(\Theta)}{R}\,\bet,\qquad
   \frac{\partial\bx}{\partial Z}=\be_{3},
\]
so
\[
   \boxed{\;\bF=f'(R)\,\ber\otimes\hbE_{R}
   +\frac{f(R)g'(\Theta)}{R}\,\bet\otimes\hbE_{\Theta}
   +\be_{3}\otimes\hbE_{3} .\;}
\]

\textbf{(b)} Both frames are orthonormal, so $\det\bF$ is the product of
the three coefficients:
\[
   \det\bF=f'\,\frac{fg'}{R}\cdot1=1 .
\]
The variables separate. Rearranged, $R/(ff')=g'(\Theta)$, and the left
side depends on $R$ alone while the right depends on $\Theta$ alone, so
both equal a constant $k$:
\[
   g'(\Theta)=k,\qquad f(R)f'(R)=\frac{R}{k} .
\]
The gap closes when the material of angular extent $\gamma$ covers
$2\pi$, so $g(0)=0$ and $g(\gamma)=2\pi$, giving
\[
   k=\frac{2\pi}{\gamma},\qquad g(\Theta)=\frac{2\pi}{\gamma}\Theta .
\]
Integrating the radial equation, $\tfrac12f^{2}=R^{2}/(2k)+\text{const}$,
and the axis must map to the axis, $f(0)=0$, which kills the constant:
\[
   \boxed{\;f(R)=\frac{R}{\sqrt k}=R\sqrt{\frac{\gamma}{2\pi}},
   \qquad g(\Theta)=\frac{2\pi}{\gamma}\Theta .\;}
\]
Hence the deformed radius is
\[
   a=f(A)=A\sqrt{\frac{\gamma}{2\pi}}\ \propto\ \sqrt\gamma ,
\]
so closing a larger gap, that is starting from a smaller $\gamma$, gives
a thinner bar. The area is preserved: $\pi a^{2}=\pi A^{2}\gamma/(2\pi)
=\tfrac12 A^{2}\gamma$, which is exactly the area of the original
sector.

\textbf{(c)} With $f'=1/\sqrt k$ and $fg'/R=(R/\sqrt k)k/R=\sqrt k$,
\[
   \bF=\frac{1}{\sqrt k}\,\ber\otimes\hbE_{R}
   +\sqrt k\,\bet\otimes\hbE_{\Theta}+\be_{3}\otimes\hbE_{3},
\]
which is already the canonical form $\bF=\sum_{i}\lambda_{i}
\bv_{i}\otimes\bu_{i}$ of \eqr{3.43}, both triads being orthonormal.
Reading off the polar factors,
\[
   \boxed{\;
\begin{aligned}
   \bR&=\ber\otimes\hbE_{R}+\bet\otimes\hbE_{\Theta}
        +\be_{3}\otimes\hbE_{3},\\
   \bU&=\frac{1}{\sqrt k}\hbE_{R}\otimes\hbE_{R}
        +\sqrt k\,\hbE_{\Theta}\otimes\hbE_{\Theta}
        +\hbE_{3}\otimes\hbE_{3},\\
   \bV&=\frac{1}{\sqrt k}\ber\otimes\ber
        +\sqrt k\,\bet\otimes\bet+\be_{3}\otimes\be_{3} .
\end{aligned}\;}
\]
Note that $\bR$ is not the shifter: it carries $\hbE_{R}(\Theta)$ to
$\ber(g(\Theta))$, and those are different directions unless $k=1$.

\textbf{(d)} The principal stretches are the coefficients:
\[
   \boxed{\;\lambda_{1}=\frac{1}{\sqrt k}=\sqrt{\frac{\gamma}{2\pi}},
   \qquad \lambda_{2}=\sqrt k=\sqrt{\frac{2\pi}{\gamma}},
   \qquad\lambda_{3}=1 .\;}
\]
Their product is $1$, as isochoricity requires. For $\gamma<2\pi$ the
material is stretched azimuthally and compressed radially by reciprocal
factors, which is a state of pure shear in the cross-sectional plane.

\textbf{(e)} With $\bB=\bV^{2}$ from (c),
\[
   T_{rr}=-p+\frac{G}{k},\qquad T_{\theta\theta}=-p+Gk,
   \qquad T_{zz}=-p+G ,
\]
off-diagonal components zero, and all quantities functions of $r$ alone.
The radial equilibrium equation is
\[
   \frac{\partial T_{rr}}{\partial r}
   +\frac{T_{rr}-T_{\theta\theta}}{r}=0,
   \qquad T_{rr}-T_{\theta\theta}=G\Big(\frac1k-k\Big),
\]
so
\[
   p'(r)=\frac{G}{r}\Big(\frac1k-k\Big),
   \qquad
   p(r)=G\Big(\frac1k-k\Big)\ln r+\text{const}.
\]
Since $\gamma<2\pi$ gives $k>1$, the coefficient $1/k-k$ is negative and
$p\to+\infty$ as $r\to0$:
\[
   \boxed{\;p\ \text{is logarithmically singular on the centreline.}\;}
\]

\emph{The outer traction.} At $r=a$ the normal is $\ber$ and
$\bT\ber=T_{rr}(a)\ber$, a normal traction, uniform because it depends
only on $r$. The constant of integration is not fixed by the field
equations, so $T_{rr}(a)$ may be given any value: the deformation is
maintained by an \emph{arbitrary uniform pressure} on the outer
boundary. This is the indeterminacy that Problem 10 removed by assuming
a traction-free patch.

\emph{The seam.} Fix the constant by requiring no traction at the outer
boundary, $T_{rr}(a)=0$, that is $p(a)=G/k$. Then
\[
   p(r)=\frac{G}{k}+G\Big(\frac1k-k\Big)\ln\frac{r}{a} .
\]
The two faces of the wedge meet on a half plane whose normal is $\bet$,
and the traction there is $\bT\bet=T_{\theta\theta}\bet$, purely normal
since $T_{r\theta}=T_{z\theta}=0$. Substituting,
\[
   \boxed{\;T_{\theta\theta}(r)
   =G\Big(k-\frac1k\Big)\Big[1+\ln\frac{r}{a}\Big],
   \qquad k=\frac{2\pi}{\gamma} .\;}
\]
So the glue transmits a pure normal traction, with no shear. It is
tensile at the rim, where $\ln(r/a)=0$ and $T_{\theta\theta}
=G(k-1/k)>0$, changes sign at $r=a/e$, and is compressive and
logarithmically unbounded towards the axis. The bar is trying to spring
open at its surface and is crushed together at its core, which is what
one would expect of a wedge forced shut.
\end{solution}

\begin{problem}[The referential energy balance for an elastic solid]
Derive the energy balance equation
\[
   \rho_{\kappa}r=\Dvg\bq_{\kappa}
   -\theta\big(W_{\theta}\big)^{\textstyle\cdot},
   \qquad\text{where}\quad
   \big(W_{\theta}\big)^{\textstyle\cdot}
   =W_{\theta\theta}\dot\theta+W_{\theta\bF}\cdot\dot\bF ,
\]
with $\bq_{\kappa}$ the referential heat flux of \eqr{6.174}.
\end{problem}

\begin{solution}
\begin{remark}[the overdot in the statement]\label{rem:ch11dotW}
The book writes the second factor as $\dot W$, but the expression given
for it, $W_{\theta\theta}\dot\theta+W_{\theta\bF}\cdot\dot\bF$, is the
material derivative of $W_{\theta}$, not of $W$. The chain rule for $W$
itself would be $\dot W=W_{\theta}\dot\theta+W_{\bF}\cdot\dot\bF$, with
first derivatives. These notes write
$(W_{\theta})^{\textstyle\cdot}$, and the derivation below shows it is
the right object.
\end{remark}

\textbf{Step 1: the internal energy in terms of $W$.} By \eqr{6.186},
$\psi=\varepsilon-\theta\eta$, so $\varepsilon=\psi+\theta\eta$.
Multiply by $\rho_{\kappa}$ and use \eqr{11.17} on the first term:
\[
   \rho_{\kappa}\varepsilon=\rho_{\kappa}\psi+\theta\rho_{\kappa}\eta
   =W+\theta\,\rho_{\kappa}\eta .
\]
The entropy is given by \eqr{11.16}$_{1}$, $\eta=-\Psi_{\theta}$, and
$\rho_{\kappa}$ passes through the $\theta$ derivative as a constant,
so $\rho_{\kappa}\eta=-\rho_{\kappa}\Psi_{\theta}=-W_{\theta}$. Hence
\[
   \rho_{\kappa}\varepsilon=W-\theta W_{\theta} .
\]

\textbf{Step 2: differentiate.} Taking the material derivative, and
using the chain rule $\dot W=W_{\theta}\dot\theta+W_{\bF}\cdot\dot\bF$
on the first term and the product rule on the second,
\[
\begin{aligned}
   \rho_{\kappa}\dot\varepsilon
   &=\dot W-\dot\theta\,W_{\theta}
     -\theta\big(W_{\theta}\big)^{\textstyle\cdot}\\
   &=\big(W_{\theta}\dot\theta+W_{\bF}\cdot\dot\bF\big)
     -W_{\theta}\dot\theta-\theta\big(W_{\theta}\big)^{\textstyle\cdot}
   =W_{\bF}\cdot\dot\bF
     -\theta\big(W_{\theta}\big)^{\textstyle\cdot} ,
\end{aligned}
\]
the two $W_{\theta}\dot\theta$ terms cancelling. Since $\bP=W_{\bF}$ by
\eqr{11.16}$_{2}$,
\[
   \rho_{\kappa}\dot\varepsilon
   =\bP\cdot\dot\bF-\theta\big(W_{\theta}\big)^{\textstyle\cdot} .
\]

\textbf{Step 3: compare with the balance.} The referential energy
balance is \eqr{6.174},
\[
   \rho_{\kappa}\dot\varepsilon
   =\bP\cdot\dot\bF-\Dvg\bq_{\kappa}+\rho_{\kappa}r .
\]
Subtracting the two displays, the stress power cancels and
\[
   -\theta\big(W_{\theta}\big)^{\textstyle\cdot}
   =-\Dvg\bq_{\kappa}+\rho_{\kappa}r ,
\]
that is
\[
   \boxed{\;\rho_{\kappa}r=\Dvg\bq_{\kappa}
   -\theta\big(W_{\theta}\big)^{\textstyle\cdot}\;}
\]
Finally, $W_{\theta}$ is a function of $\theta$ and $\bF$, so the chain
rule gives $(W_{\theta})^{\textstyle\cdot}=W_{\theta\theta}\dot\theta
+W_{\theta\bF}\cdot\dot\bF$, as stated.

\textbf{Reading the result.} The stress power has disappeared entirely.
In a hyperelastic solid the mechanical working is stored, not dissipated,
so it cannot appear as a source of heat; what remains is the conduction
term and a coupling term $-\theta(W_{\theta})^{\textstyle\cdot}$,
which is the rate of latent heat associated with changes of temperature
and deformation. Its second piece, $-\theta W_{\theta\bF}\cdot\dot\bF$,
is the thermoelastic coupling: deforming the body at fixed temperature
absorbs or releases heat unless $W_{\theta\bF}=\hOb$, in which case the
thermal and mechanical problems decouple as they did for rigid bodies at
Section~\ref{sec:ch6energy}.
\end{solution}

\begin{problem}[The referential heat flux]
Show that if the spatial heat flux is \eqr{11.49}, then the referential
heat flux is
\[
   \bq_{\kappa}=\big(\psi_{0}\hI+\psi_{1}\bC+\psi_{2}\bC^{2}\big)
   \Grad\theta ,
\]
with $\psi_{0},\psi_{1},\psi_{2}$ functions of $\theta$ and of the
invariants
\[
   I_{1}(\bC),\quad I_{2}(\bC),\quad I_{3}(\bC),\quad
   I_{4}(\bC,\Grad\theta),\quad I_{5}(\bC,\Grad\theta),\quad
   I_{6}(\Grad\theta).
\]
\end{problem}

\begin{solution}
\textbf{Two conversions.} By \eqr{6.171}, $\bq_{\kappa}=J\bF^{-1}\bq$.
By the chain rule applied to the temperature field,
$\Grad\theta=\bF^{t}\grad\theta$, whence
$\grad\theta=\bF^{-t}\Grad\theta$. Substituting \eqr{11.49},
\[
   \bq_{\kappa}=J\bF^{-1}
   \big(\phi_{0}\I+\phi_{1}\bB+\phi_{2}\bB^{2}\big)
   \bF^{-t}\,\Grad\theta .
\]

\textbf{Pulling the three tensors back.} Each is a one-line
computation, the reverse of the ones in Problem 2:
\[
   \bF^{-1}\I\bF^{-t}=\bF^{-1}\bF^{-t}=\big(\bF^{t}\bF\big)^{-1}
   =\bC^{-1},
\]
\[
   \bF^{-1}\bB\bF^{-t}=\bF^{-1}\big(\bF\bF^{t}\big)\bF^{-t}=\hI,
\]
\[
   \bF^{-1}\bB^{2}\bF^{-t}
   =\bF^{-1}\bF\bF^{t}\bF\bF^{t}\bF^{-t}=\bF^{t}\bF=\bC .
\]
Therefore
\[
   \bq_{\kappa}=J\big(\phi_{0}\bC^{-1}+\phi_{1}\hI
   +\phi_{2}\bC\big)\Grad\theta .
\]

\textbf{Removing the inverse.} The result is not yet in the stated form,
because of $\bC^{-1}$. The Cayley--Hamilton theorem of
Section~\ref{sec:ch9isotropic}, applied to $\bC$, gives
$\bC^{3}-I_{1}\bC^{2}+I_{2}\bC-I_{3}\hI=\hOb$, and multiplying by
$\bC^{-1}$ and rearranging,
\[
   \bC^{-1}=\frac{1}{I_{3}}\big(\bC^{2}-I_{1}\bC+I_{2}\hI\big),
\]
legitimate because $I_{3}=\det\bC=J^{2}>0$. Substituting and collecting
powers of $\bC$,
\[
   \boxed{\;\bq_{\kappa}=\big(\psi_{0}\hI+\psi_{1}\bC
   +\psi_{2}\bC^{2}\big)\Grad\theta\;}
\]
with
\[
   \psi_{0}=J\Big(\phi_{1}+\frac{I_{2}}{I_{3}}\phi_{0}\Big),
   \qquad
   \psi_{1}=J\Big(\phi_{2}-\frac{I_{1}}{I_{3}}\phi_{0}\Big),
   \qquad
   \psi_{2}=\frac{J}{I_{3}}\phi_{0} .
\]

\textbf{The argument lists match.} The $\phi_{k}$ depend on $\theta$, on the invariants of $\bB$ and on
the three scalars \eqr{11.50}. The invariants of $\bB$ are those of
$\bC$, by \eqr{11.35}. For the others, substitute
$\grad\theta=\bF^{-t}\Grad\theta$ and pull back as above:
\[
\begin{aligned}
   I_{4}(\bB,\grad\theta)
   &=\ip{\Grad\theta}{\bF^{-1}\bB\bF^{-t}\Grad\theta}
   =\ip{\Grad\theta}{\Grad\theta}=I_{6}(\Grad\theta)^{2},\\
   I_{5}(\bB,\grad\theta)
   &=\ip{\Grad\theta}{\bC\,\Grad\theta}=I_{4}(\bC,\Grad\theta),\\
   I_{6}(\grad\theta)^{2}
   &=\ip{\Grad\theta}{\bC^{-1}\Grad\theta}\\
   &=\frac{1}{I_{3}}\big[I_{5}(\bC,\Grad\theta)
   -I_{1}I_{4}(\bC,\Grad\theta)+I_{2}I_{6}(\Grad\theta)^{2}\big],
\end{aligned}
\]
the last line using the same Cayley--Hamilton substitution, and with the
referential scalars defined by
$I_{4}(\bC,\Grad\theta)=\ip{\Grad\theta}{\bC\Grad\theta}$,
$I_{5}(\bC,\Grad\theta)=\ip{\Grad\theta}{\bC^{2}\Grad\theta}$ and
$I_{6}(\Grad\theta)=|\Grad\theta|$. Every argument of the $\phi_{k}$ is
therefore a function of the referential list, and so is every
$\psi_{k}$.

The moral is the one of Remark~\ref{rem:ch6threestress} transposed to
fluxes: the referential and spatial descriptions carry the same physical
content, and converting between them replaces $\bB$ by $\bC$ and turns
$\grad$ into $\Grad$, at the price of factors of $\bF$ that
Cayley--Hamilton then tidies away.
\end{solution}

\begin{problem}[A torsion with an axial temperature gradient]
A solid circular cylindrical bar undergoes the static torsion of
Section~\ref{sec:torsion}. The $\phi$'s in \eqr{11.49} are independent
of temperature, though they may depend on $I_{1}$ to $I_{6}$, and the
material is uniform. The temperature is independent of time and depends
only on $z$, with the ends $z=0$ and $z=l$ held at $\theta_{0}$ and
$\theta_{l}$. Assuming this function linear, show that the energy
balance with zero internal heating supply is automatically satisfied,
and that the lateral surface of the bar is insulated.
\end{problem}

\begin{solution}
\textbf{The data.} From Problem 12,
\[
   \bB=\I+\gamma\big(\bet\otimes\be_{3}+\be_{3}\otimes\bet\big)
   +\gamma^{2}\bet\otimes\bet,\qquad \gamma=\tau r ,
\]
and the temperature is
\[
   \theta(z)=\theta_{0}+\beta z,\qquad
   \beta=\frac{\theta_{l}-\theta_{0}}{l},\qquad
   \grad\theta=\beta\,\be_{3},
\]
a constant vector field.

\textbf{The heat flux.} Apply $\bB$ and $\bB^{2}$ to $\be_{3}$:
\[
   \bB\be_{3}=\be_{3}+\gamma\bet,\qquad
   \bB\bet=(1+\gamma^{2})\bet+\gamma\be_{3},
\]
so
\[
\begin{aligned}
   \bB^{2}\be_{3}&=\bB\big(\be_{3}+\gamma\bet\big)
   =\big(\be_{3}+\gamma\bet\big)
   +\gamma\big[(1+\gamma^{2})\bet+\gamma\be_{3}\big]\\
   &=(1+\gamma^{2})\be_{3}+\gamma(2+\gamma^{2})\bet .
\end{aligned}
\]
Substituting into \eqr{11.49},
\[
   \boxed{\;\bq=\beta\Big\{\big[\phi_{0}+\phi_{1}
   +\phi_{2}(1+\gamma^{2})\big]\be_{3}
   +\gamma\big[\phi_{1}+\phi_{2}(2+\gamma^{2})\big]\bet\Big\}\;}
\]

\textbf{The lateral surface is insulated.} Its outward normal is
$\ber$, and the flux has no $\ber$ component:
\[
   \ip\bq\ber=0 ,
\]
so no heat crosses the cylindrical surface. The flux runs along the
axis, as the boundary data demand, and around it, which is the
thermal signature of the torsion.

\textbf{The energy balance.} The local spatial form is \eqr{6.168},
\[
   \rho\dot\varepsilon=\bT\cdot\bL-\dvg\bq+\rho r ,
\]
with $r=0$ by hypothesis. The deformation is static, so $\bv=\bze$,
$\bL=\Ob$ and the stress power vanishes; and $\varepsilon$ depends on
$\theta$ and $\bF$, neither of which depends on time, so
$\dot\varepsilon=0$. What is left is $\dvg\bq=0$, and this holds
identically. To see it, note first what the coefficients depend on. The
$\phi_{k}$ are independent of $\theta$ by hypothesis and of $\bX$ by
uniformity, so they depend only on the six invariants, and
\[
\begin{aligned}
   I_{1}=I_{2}&=3+\gamma^{2},\qquad I_{3}=1,\\
   I_{4}&=\ip{\grad\theta}{\bB\grad\theta}=\beta^{2}
   \ip{\be_{3}}{\be_{3}+\gamma\bet}=\beta^{2},\\
   I_{5}&=\ip{\grad\theta}{\bB^{2}\grad\theta}
   =\beta^{2}(1+\gamma^{2}),\qquad I_{6}=|\beta| ,
\end{aligned}
\]
the first line by \eqr{11.38}, torsion being locally a simple shear of
amount $\gamma$. Every one of these is a function of $r$ alone, through
$\gamma=\tau r$. So the boxed flux has the form
\[
   \bq=q_{\theta}(r)\,\bet+q_{z}(r)\,\be_{3},\qquad q_{r}=0 ,
\]
and the cylindrical divergence of such a field is
\[
   \dvg\bq=\frac{1}{r}\frac{\partial}{\partial r}\big(r\,q_{r}\big)
   +\frac{1}{r}\frac{\partial q_{\theta}}{\partial\theta}
   +\frac{\partial q_{z}}{\partial z}
   =0+0+0=0 ,
\]
the first term because $q_{r}=0$, the second because $q_{\theta}$ does
not depend on $\theta$, the third because $q_{z}$ does not depend on
$z$. Hence
\[
   \boxed{\;\dvg\bq=0\ \text{ identically, and the balance is
   satisfied with }r=0 .\;}
\]

\textbf{Why it worked.} Three separate facts conspired. The temperature
gradient is uniform, so it contributes nothing to any derivative. The
deformation is a function of $r$ only, so the constitutive coefficients
are too. And the flux has no radial component, so the one term of the
divergence that would have differentiated an $r$-dependent quantity is
absent. Remove any one of them, for instance by letting the $\phi$'s
depend on temperature, and $\dvg\bq$ would acquire a term in
$\partial_{z}$ through $\theta(z)$, and an internal supply $r$ would be
needed to maintain the assumed field.
\end{solution}
